\documentclass[a4paper,12pt]{article}
\usepackage[utf8]{inputenc}
\usepackage[german]{babel}
\usepackage{amsmath}
\usepackage{amssymb}
\usepackage{amsthm}
\usepackage{geometry}
\usepackage{graphicx}
\usepackage[hidelinks]{hyperref}
\usepackage{listings}
\usepackage{xcolor}
\usepackage{algorithm}
\usepackage{algpseudocode}
\usepackage{chngcntr}
\usepackage{microtype}
\usepackage{enumitem}
\setlist[itemize,1]{label=-}

\geometry{margin=2.5cm}

\hypersetup{
	colorlinks=true,
	linkcolor=blue!60!black,
	urlcolor=blue!60!black,
	citecolor=blue!60!black
}

\theoremstyle{definition}
\newtheorem{definition}{Definition}
\newtheorem{beispiel}{Beispiel}
\newtheorem{frage}{Frage}
\newtheorem{satz}{Satz}
\newtheorem{lemma}{Lemma}
\newtheorem{korollar}{Korollar}
\counterwithin{algorithm}{section}

% Code-Highlighting
\lstset{
	basicstyle=\ttfamily\small,
	keywordstyle=\color{blue},
	commentstyle=\color{gray},
	stringstyle=\color{red},
	numbers=left,
	numberstyle=\tiny,
	breaklines=true,
	frame=single
}

\title{Graphen mit Knoten und Kanten \\[0.5ex]\large Optimale Einordnung in die Graphprofilverteilung}
\author{Stephan Epp}
\date{\today}

\begin{document}
	
	\maketitle
	
	\tableofcontents
	\newpage
	
	\section{Einführung}

	\subsection{Graphen als universelles Modell}

	Graphen gehören zu den mächtigsten und vielseitigsten Abstraktionen der Informatik und Mathematik. Ein Graph besteht aus einer Menge von Knoten und einer Menge von Kanten, die Paare von Knoten verbinden. Diese schlichte Struktur genügt, um eine außerordentlich breite Klasse realer Systeme zu modellieren: Straßennetze, Computernetzwerke, soziale Beziehungsgeflechte, Proteininteraktionen, logische Abhängigkeiten in Software sowie – besonders eindrucksvoll – die synaptischen Verbindungen zwischen Neuronen im menschlichen Gehirn.

	Die Stärke des Graphmodells liegt in seiner Abstraktion. Sobald ein System als Graph formalisiert ist, lassen sich allgemeine algorithmische Werkzeuge anwenden, unabhängig vom jeweiligen Anwendungsgebiet. Fragen wie \emph{Wie weit sind zwei Knoten voneinander entfernt?}, \emph{Gibt es einen Weg von $A$ nach $B$?} oder \emph{Welche Knoten sind für die Konnektivität des Netzes besonders kritisch?} übersetzen sich direkt in algorithmische Probleme auf Graphen, deren Lösung in allen Anwendungsdomänen gleichermaßen nützlich ist.

	\subsection{Das Problem der strukturellen Charakterisierung}

	Wissen, dass ein System als Graph modelliert werden kann, reicht allein nicht aus. Um das Verhalten eines Netzwerks zu verstehen, es mit anderen Netzwerken zu vergleichen oder Entscheidungen über seine Gestaltung zu treffen, benötigt man quantitative Kenngrößen, die die \emph{Struktur} des Graphen präzise beschreiben. Hierfür sind insbesondere zwei Eigenschaften von zentraler Bedeutung:

	\begin{itemize}
		\item[-] \textbf{Weglängen}: Der kürzeste Weg zwischen zwei Knoten $i$ und $j$ gibt an, wie viele Schritte mindestens nötig sind, um von $i$ nach $j$ zu gelangen. Der längste (einfache) Weg beschreibt hingegen die maximale Tiefe erreichbarer Kettenstrukturen. Beide Maße zusammen charakterisieren den \emph{Informationsfluss} in einem Netzwerk.
		\item[-] \textbf{Kantendichte}: Das Verhältnis von Knoten- zu Kantenzahl, hier als Kantenmaß $\kappa = |V| / |E|$ definiert, quantifiziert, wie \emph{sparsam} oder \emph{dicht} ein Graph verdrahtet ist. Ein kleines $\kappa$ bedeutet viele Verbindungen pro Knoten; ein großes $\kappa$ kennzeichnet spärlich verbundene, ressourcenschonende Topologien.
	\end{itemize}

	Das \emph{Profil} eines Graphen, definiert als das Tripel aus der kürzeste-Wege-Matrix $D$, der längste-Wege-Matrix $L$ und dem Kantenmaß $\kappa$, fasst diese Eigenschaften zu einer einzigen, kohärenten Beschreibung zusammen. Ziel dieser Arbeit ist es zu zeigen, dass sich dieses Profil effizient und \emph{optimal} berechnen lässt und wie es zur Einordnung von Graphen in eine \emph{Graphprofilverteilung} dient.

	\subsection{Rechenaufwand und die Frage nach effizienten Algorithmen}

	Die naive Berechnung aller kürzesten Wege in einem Graphen mit $n$ Knoten erfordert klassisch $O(n^3)$ Operationen, beispielsweise durch den Algorithmus von \textsc{Floyd-Warshall} oder durch $n$-fache Anwendung von \textsc{Dijkstra} auf dichten Graphen. Für kleine bis mittlere Graphen ($n \leq 10{.}000$) ist dies auf moderner Hardware in Sekunden durchführbar; für sehr große Netzwerke – etwa mit $n = 10^6$ oder gar $n = 10^9$ Knoten – werden diese Methoden unpraktikabel.

	Einen entscheidenden Schlüssel zu einer schnelleren Variante liefert die \emph{Boolean Matrixmultiplikation}. Anstatt numerische Gewichte zu akkumulieren, genügt es für viele graphentheoretische Fragestellungen, lediglich zu entscheiden, ob ein Weg \emph{existiert} oder nicht – eine Ja-/Nein-Information, die sich mit logischen Operationen ($\wedge$, $\vee$) verarbeiten lässt. Wiederholt man die Boolean Matrixmultiplikation für wachsende Potenzen der Adjazenzmatrix $A$, ergibt die $k$-te Potenz $A^k$ genau an, welche Knotenpaare durch Wege der Länge $k$ verbunden sind.

	Die entscheidende Beobachtung dieser Arbeit ist, dass sich jede Zeile von $A$ als ganzzahlige \emph{Signatur} kodieren lässt – eine Binärzahl, bei der das $k$-te Bit gesetzt ist, genau wenn der $k$-te Eintrag der Zeile~$1$ ist. Die Boolean Matrixmultiplikation zweier Matrizen reduziert sich dann auf die bitweise UND-Operation zweier solcher Signaturen, was auf moderner Hardware in $O(1)$ abläuft. Diese \emph{Signatur-Methode} führt zu einer Boolean Matrixmultiplikation in $O(n^2)$ – gegenüber $O(n^3)$ bei naiver Implementierung eine Verbesserung um einen vollen Faktor $n$.

	\subsection{Determinismus als algorithmisches Gütekriterium}

	In der Theoretischen Informatik unterscheidet man deterministische von nichtdeterministischen Berechnungsmodellen. Einem deterministischen Algorithmus ist für jede Eingabe genau ein Berechnungspfad zugeordnet; das Ergebnis ist reproduzierbar und vorhersagbar. Nichtdeterministische oder probabilistische Algorithmen hingegen erzeugen je nach Zufallskomponente möglicherweise unterschiedliche Resultate auf derselben Eingabe.

	Für die Praxis bedeutet dies: \emph{Deterministische Algorithmen liefern für alle Instanzen eine Garantie}, während probabilistische Verfahren nur mit einer gewissen Wahrscheinlichkeit korrekte Ergebnisse liefern. In sicherheitskritischen Anwendungen – von der medizinischen Bildgebung über die Infrastrukturplanung bis hin zur Verifikation von KI-Systemen – ist diese Garantie unverzichtbar. Die in dieser Arbeit entwickelte Methode zur Graphprofilberechnung ist vollständig deterministisch: Sie berechnet für jeden Graphen $G$ stets exakt dasselbe Profil $(D, L, \kappa)$, ohne jede Zufallskomponente.

	\subsection{Biologische Motivation: Das menschliche Gehirn als Graph}

	Eine der faszinierendsten Anwendungen der Graphtheorie ist die Analyse neuronaler Netzwerke. Das menschliche Gehirn besteht aus schätzungsweise $86 \cdot 10^9$ Neuronen, die durch ca. $1{,}5 \cdot 10^{14}$ synaptische Verbindungen miteinander verknüpft sind – ein gewaltiger Graph, dessen strukturelle Eigenschaften unmittelbar mit kognitiven Fähigkeiten zusammenhängen. Die Frage, welches \emph{Profil} neuronale Netzwerke aufweisen und wie dieses Profil mit biochemischen Kosten und Verarbeitungsgeschwindigkeit zusammenhängt, ist Gegenstand intensiver Forschung in den Neurowissenschaften.

	Jede synaptische Verbindung – jede Kante dieses gigantischen Graphen – ist kein passiver Kanal, sondern ein dynamischer Übertragungsort des \emph{unsichtbaren Geistes des Lebens}: des elektro-chemischen Informationsflusses von Aktionspotential zu Aktionspotential. Ein einzelnes Aktionspotential dauert $\approx 1$\,ms und erzeugt ein Membranpotential von ca. $+40$\,mV – für das bloße Auge unsichtbar, aber die Trägersubstanz jedes Gedankens, jeder Erinnerung, jeder Entscheidung. Die synaptische Stärke – das Kantengewicht im neuronalen Graphen – variiert durch Langzeitpotenzierung (LTP, verstärkte Verbindungen) und Langzeitdepression (LTD, abgeschwächte Verbindungen) sowie durch Spike-Timing Dependent Plasticity (STDP). Das Gehirn ist damit ein \emph{gewichteter, plastischer Graph}, dessen Kantengewichte sich durch Erfahrung stetig verändern. Die mathematische Grundlage beschreiben die Hodgkin-Huxley-Gleichungen, die für jeden Knoten (Neuron) die zeitliche Entwicklung des Membranpotentials $V(t)$ als Funktion der eingehenden Ionenströme berechenbar machen.

	Eine graphentheoretisch bedeutsame Beobachtung ist die \emph{Rechtsdrehung kortikaler Informationsverarbeitung}: Von dorsal (oben) betrachtet propagiert die kortikale Aktivierungswelle im Uhrzeigersinn – von anterior über lateral-rechts zu posterior, dann lateral-links und zurück nach anterior. Formal entspricht dies einer negativen Winkelgeschwindigkeit des kortikalen Aktivierungsvektors $\vec{A}(t)$:
	\[
	\omega(t) \;=\; \frac{A_x(t)\,\dot{A}_y(t) - A_y(t)\,\dot{A}_x(t)}{|\vec{A}(t)|^2} \;<\; 0.
	\]
	Das gesunde Gehirn ist damit ein \emph{chiraler, gerichteter Graph}: Der Informationsfluss besitzt eine strukturell bevorzugte Richtung, und ein vollständiger kortikaler Zyklus dauert bei $\approx 40$\,Hz nur ca. $25$\,ms – genau ein Gamma-Fenster. Anatomisch liegt die Ursache in der Asymmetrie thalamo-kortikaler Schleifen (höhere Projektionsdichte im rechten Thalamus), hämodynamisch in der leichten Rechtspräferenz des Blutflusses durch den Circulus Willisii.

	Gerät dieser Informationsfluss aus dem Gleichgewicht – durch Informationsmangel, chronischen Stress oder erhöhte Cortisolspiegel –, entsteht die \emph{epistemische Wolke}: ein messbarer Zustand verminderter kognitiver Kapazität. Die Wolken-Dichte
	\[
	\delta(\mathcal{W}) \;:=\; 1 - \frac{|\mathcal{I}_P \cap \mathcal{I}^*|}{|\mathcal{I}^*|} \;\in\; [0,1]
	\]
	quantifiziert die Abweichung der verfügbaren Informationsmenge $\mathcal{I}_P$ vom notwendigen Set $\mathcal{I}^*$. Für $\delta = 0$ herrscht vollständige Klarheit; für $\delta = 1$ vollständige Blockade. Graphentheoretisch manifestiert sich die epistemische Wolke als dreifache Störung: erhöhter Netzwerkdurchmesser (längere Reaktionszeiten beim Informationsabruf), reduzierte effektive Kantenkapazität (EEG-Gamma-Kohärenz sinkt um $\approx 60\,\%$) und destabilisierte chiral-rechtsdrehende Traversierungsrichtung. Neurobiologische Signaturen sind erhöhte Amygdala-Reaktivität, reduzierte präfrontale Konnektivität, supprimierte Gamma-Kohärenz und chronisch erhöhte Cortisolspiegel, die ihrerseits synaptische Dichte und Neurotransmitter-Biosynthese weiter beeinträchtigen.

	Eine fünfte fundamentale Dimension ist die \emph{kognitive Benetzung} $\mathcal{B}$ – der Grad, zu dem alle $1{,}5 \cdot 10^{14}$ synaptischen Verbindungen gleichmäßig mit aktivem Informationsfluss durchdrungen sind. Analog zur Perkolationstheorie gilt: Nur wenn $\mathcal{B} \geq \mathcal{B}_{\min} \approx 0{,}75$ (Aktivierungswahrscheinlichkeit jeder Synapse oberhalb des Perkolationsschwellwerts), existiert im neuronalen Graphen eine giant component der Größe $\Theta(N)$ – die Grundlage unbegrenzten freien Denkens. Unterhalb dieses Schwellwerts kollabiert die kognitive Kapazität phasenübergangsartig. Die epistemische Wolke wirkt als kognitives Trocknungsmittel: Cortisol hemmt Tyrosin-Hydroxylase (Schlüsselenzym der Dopaminsynthese) und Tryptophan-Hydroxylase (Serotoninsynthese); erhöhte GABA-Aktivität hyperpolarisiert Neurone unter die Feuerschwelle; Astrozyten-Dysfunktion fragmentiert das Netzwerk in nicht kommunizierende Inseln. Das vollständige Modell freien Denkens lautet:
	\[
	\Psi_{\mathrm{frei}}(t) \;=\; \mathcal{B}(\delta,\,T,\,V_m) \;\cdot\; \mathrm{IZI}(t) \;\cdot\; (1-\delta) \;\cdot\; \mathbf{1}[\omega(t) < 0],
	\]
	wobei $\mathrm{IZI}$ der Informationsfluss-Intensitäts-Index (synaptischer Lebensgeist-Fluss) und $\mathbf{1}[\omega < 0]$ der Rechtsdrehungs-Indikator ist. Alle vier Faktoren sind multiplikativ verknüpft: Fällt einer auf null, kollabiert $\Psi_{\mathrm{frei}}$ vollständig.

	Im Rahmen dieser Arbeit werden sechs biologisch motivierte Graphklassen untersucht: zufällige Sparse-Graphen (\textsc{Erdős-Rényi}), skalenfreie Netzwerke (\textsc{Barabási-Albert}), Small-World-Graphen (\textsc{Watts-Strogatz}), hierarchische Netzwerke, modulare Netzwerke und kortikale Säulenstrukturen. Für jede Klasse wird das Profil systematisch berechnet und ausgewertet. Die Ergebnisse erlauben es, reale Gehirnnetzwerke in die Graphprofilverteilung einzuordnen und charakteristische Muster biologischer Netzwerkarchitekturen zu identifizieren.

	Auf diesem Weg zeigt die Arbeit exemplarisch, dass das Graphprofil nicht nur ein theoretisches Konstrukt ist, sondern ein \emph{praktisch einsetzbares Werkzeug}: Die experimentellen Daten belegen einen robusten inversen Zusammenhang zwischen Kantenmaß $\kappa$ und Netzwerkdurchmesser, der sich durch eine einfache Regressionsformel beschreiben lässt und direkte Implikationen für das Design effizienter neuronaler Architekturen hat.

	\subsection{Beiträge und Aufbau dieser Arbeit}

	Die zentralen Beiträge dieser Arbeit lassen sich wie folgt zusammenfassen:

	\begin{enumerate}
		\item \textbf{Effiziente Profilberechnung}: Es wird gezeigt, wie die Signatur-Methode aus der Boolean Matrixmultiplikation zur simultanen Berechnung von kürzesten Wegen, längsten Wegen und Kantenmaß in einem einheitlichen Algorithmus (Algorithmus~\ref{alg:profile}) genutzt werden kann. Die Gesamtlaufzeit beträgt $O(n^3)$ bei $O(n^2)$ Speicherbedarf.
		\item \textbf{Optimale Einordnung}: Es wird bewiesen, dass die so berechnete Einordnung jedes Graphen in die Graphprofilverteilung optimal und nicht verbesserbar ist – jede auf dem Profil basierende strukturelle Aussage ist maximal informiert.
		\item \textbf{Experimentelle Validierung}: An sechs biologisch inspirierten Graphklassen wird die theoretische Laufzeitanalyse experimentell bestätigt (gemessener Exponent $\approx n^{3.16}$). Die charakteristischen Profile biologischer Netzwerke werden quantifiziert und in die Graphprofilverteilung eingeordnet.
		\item \textbf{Anwendungsperspektiven}: Konkrete Einsatzmöglichkeiten in den Neurowissenschaften, der Künstlichen Intelligenz, dem Rechenzentrumsdesign, der Sozialen Netzwerkanalyse und der Molekularbiologie werden diskutiert.
	\end{enumerate}

	\bigskip
	\noindent
	Die Arbeit ist wie folgt gegliedert: Kapitel~2 führt die grundlegenden Definitionen ein. Kapitel~3 beschreibt die Algorithmen zur Boolean Matrixmultiplikation und Profilberechnung. Kapitel~4 präsentiert die experimentellen Ergebnisse. Kapitel~5 analysiert die optimale Graphprofilverteilung und ihre theoretischen Grundlagen. Kapitel~6 diskutiert Anwendungen in verschiedenen Domänen. Kapitel~7 behandelt theoretische Implikationen, und Kapitel~8 gibt einen Ausblick auf offene Fragen. Kapitel~9 fasst die Ergebnisse zusammen.

	\section{Definitionen}
	In diesem Kapitel werden die Definitionen für diese Arbeit eingeführt.
	\subsection{Graphen}
	\begin{definition}[Graph]
		Ein Graph $G = (V, E)$ besteht aus einer Menge von Knoten $V$ und einer Menge von Kanten $E \subseteq V \times V$.
	\end{definition}
	
	\begin{definition}[Adjazenzmatrix]
		Für einen Graphen $G = (V, E)$ mit $n = |V|$ Knoten ist die Adjazenzmatrix $A \in \{0,1\}^{n \times n}$ definiert durch:
		\[
		A_{ij} = \begin{cases} 1 & \text{falls } (i,j) \in E \\ 0 & \text{sonst} \end{cases}
		\]
	\end{definition}
	
	\subsection{Wegeigenschaften}
	\begin{definition}[Weg]
		Ein Weg der Länge $k$ von Knoten $i$ zu Knoten $j$ ist eine Folge von Knoten $(v_0, v_1, \ldots, v_k)$ mit $v_0 = i$, $v_k = j$ und $(v_l, v_{l+1}) \in E$ für alle $0 \leq l < k$.
	\end{definition}
	
	\begin{lemma}[Wege und Matrixpotenzen]
		Sei $A$ die Adjazenzmatrix eines Graphen $G = (V,E)$ und $A^k$ die $k$-te Potenz von $A$ unter Boolean Matrixmultiplikation. Dann gilt:
		\[
		(A^k)_{ij} = 1 \Leftrightarrow \text{es existiert ein Weg der Länge } k \text{ von } i \text{ nach } j
		\]
	\end{lemma}
	
	\begin{proof}
		Der Beweis erfolgt durch Induktion über $k$.\newline
		\textbf{Induktionsanfang:} Für $k=1$ gilt $(A^1)_{ij} = A_{ij} = 1$ genau dann, wenn $(i,j) \in E$, also ein Weg der Länge 1 existiert.\newline
		\textbf{Induktionsschritt:} Sei die Aussage für $k$ erfüllt. Für $k+1$ gilt:
		\begin{align*}
			(A^{k+1})_{ij} &= \bigvee_{l=1}^{n} ((A^k)_{il} \wedge A_{lj}) = 1
		\end{align*}
		Dies ist genau dann erfüllt, wenn es ein $l$ gibt mit $(A^k)_{il} = 1$ und $A_{lj} = 1$. Nach Induktionsvoraussetzung existiert dann ein Weg der Länge $k$ von $i$ nach $l$ und eine Kante von $l$ nach $j$, also ein Weg der Länge $k+1$ von $i$ nach $j$.
	\end{proof}
	
	\begin{definition}[Profil]
		Das Profil eines Graphen $G=(V, E)$ besteht aus 
		\begin{itemize}
			\item[-] dem kürzesten Weg für alle Knoten $v \in V$,
			\item[-] dem längsten Weg für alle Knoten $v \in V$ und
			\item[-] dem Kantenmaß $\kappa = \frac{\left|V\right|}{\left|E\right|}$.
		\end{itemize}
	\end{definition}
	
	Mit diesem Profil wird eine Verteilung definiert, in die sich alle Graphen einordnen lassen. Ausschlaggebend für die unterschiedlichen Bereiche in der Verteilung ist die erforderliche Laufzeit der Algorithmen zur Ermittlung aller Profilwerte. Natürlich ist es dafür nötig, dass die Algorithmen die Profilwerte in optimaler Laufzeit berechnen. Sonst wäre die Einordnung nicht klar.
	
	\begin{definition}[Erreichbarkeitsmatrix]
		Die Erreichbarkeitsmatrix $R \in \{0,1\}^{n \times n}$ eines Graphen $G$ ist definiert als:
		\[
		R_{ij} = \begin{cases} 1 & \text{falls ein Weg von } i \text{ nach } j \text{ existiert} \\ 0 & \text{sonst} \end{cases}
		\]
	\end{definition}
	
	\begin{satz}[Erreichbarkeit durch Matrixpotenzen]
		Für einen Graphen mit $n$ Knoten gilt:
		\[
		R = A \vee A^2 \vee A^3 \vee \cdots \vee A^{n-1}
		\]
		wobei $\vee$ die komponentenweise logische ODER-Operation bezeichnet.
	\end{satz}
	
	\begin{proof}
		Da der längste einfache Weg in einem Graphen mit $n$ Knoten höchstens die Länge $n-1$ hat, genügt es alle Potenzen bis $A^{n-1}$ zu betrachten. Ein Weg von $i$ nach $j$ existiert genau dann, wenn es einen Weg einer Länge $k \in \{1, \ldots, n-1\}$ gibt.
	\end{proof}
	
	Die naive Implementierung hat eine Laufzeit von $O(n^3)$. Fortgeschrittene Algorithmen wie \textsc{Strassen} (1969) oder \textsc{Coppersmith-Winograd} (1990) erreichen $O(n^{2.807})$ bzw. $O(n^{2.376})$.
	
	Für \textbf{Boolean Matrizen}, bei denen alle Einträge $\{0, 1\}$ sind und die Operationen durch logische Operationen ersetzt werden, kann die Signatur-Technik aus dem Subgraph Algorithmus genutzt werden, um eine Laufzeit von $O(n^2)$ zu erreichen.
	
	\begin{definition}[Boolean Matrixmultiplikation]
		Für zwei Boolean Matrizen $A, B \in \{0,1\}^{n \times n}$ ist die Boolean Matrixmultiplikation definiert als:
		\[
		C_{ij} = \bigvee_{k=1}^{n} (A_{ik} \wedge B_{kj})
		\]
		wobei $\vee$ das logische ODER und $\wedge$ das logische UND bezeichnet.
	\end{definition}
	
	Mit anderen Worten: $C_{ij} = 1$ genau dann, wenn es mindestens ein $k$ gibt mit $A_{ik} = 1$ und $B_{kj} = 1$.
	
	\section{Algorithmen}
	
	In diesem Kapitel werden die Algorithmen zur Bestimmung der Profilwerte beschrieben.
	
	\subsection{Boolean Matrixmultiplikation}
	
	Der Kerngedanke für die Arbeitsweise des Algorithmus zur Boolean Matrixmultiplikation mit Signaturen ist, jede \textbf{Zeile} von $A$ als Signatur zu kodieren und jede \textbf{Spalte} von $B$ als Signatur zu kodieren und für jedes Element Element $C_{ij}$:
	\begin{itemize}
		\item[-] Berechne bitweise UND der Signaturen: $\sigma(\text{row}_i(A)) \; \& \; \sigma(\text{col}_j(B))$
		\item[-] Setze $C_{ij} = 1$ falls Ergebnis $\neq 0$, sonst $C_{ij} = 0$
	\end{itemize}
	
	Algorithmus \ref{alg:bmm} beschreibt die Arbeitsweise zur Boolean Matrixmultiplikation mit Signaturen.
	
	\begin{algorithm}[htb]
		\caption{Boolean Matrixmultiplikation mit Signaturen}
		\label{alg:bmm}
		\textbf{Eingabe:} Zwei Boolean Matrizen $A, B \in \{0,1\}^{n \times n}$ \\
		\textbf{Ausgabe:} Boolean Matrix $C \in \{0,1\}^{n \times n}$ mit $C_{ij} = \bigvee_{k=1}^{n} (A_{ik} \wedge B_{kj})$
		\begin{algorithmic}[1]
			\State $n \gets \text{Dimension von } A$
			\State $\text{rowSig} \gets \text{leeres Array der Länge } n$
			\State $\text{colSig} \gets \text{leeres Array der Länge } n$
			\For{$i = 0$ \textbf{to} $n-1$}
			\State $\text{rowSig}[i] \gets \sum_{k=0}^{n-1} A_{ik} \cdot 2^k$ \Comment{$O(n)$}
			\EndFor
			\For{$j = 0$ \textbf{to} $n-1$}
			\State $\text{colSig}[j] \gets \sum_{k=0}^{n-1} B_{kj} \cdot 2^k$ \Comment{$O(n)$}
			\EndFor
			\State $C \gets n \times n$ Nullmatrix
			\For{$i = 0$ \textbf{to} $n-1$}
			\For{$j = 0$ \textbf{to} $n-1$}
			\State $\text{andResult} \gets \text{rowSig}[i] \; \& \; \text{colSig}[j]$ \Comment{$O(1)$ Bitoperation}
			\If{$\text{andResult} \neq 0$}
			\State $C_{ij} \gets 1$
			\EndIf
			\EndFor
			\EndFor    
			\State \Return $C$
		\end{algorithmic}
	\end{algorithm}
	
	\begin{satz}[Korrektheit der Signatur-Methode]
		Sei $r = \sigma(\text{row}_i(A))$ die Signatur der $i$-ten Zeile von $A$ und $c = \sigma(\text{col}_j(B))$ die Signatur der $j$-ten Spalte von $B$. Dann gilt:
		\begin{align}
			C_{ij} = 1 \Leftrightarrow (r \; \& \; c) \neq 0
		\end{align}
	\end{satz}
	
	\begin{proof}
		Nach Definition ist:
		\begin{align}
			C_{ij} = 1 \Leftrightarrow \exists k: A_{ik} = 1 \text{ und } B_{kj} = 1
		\end{align}
		Die bitweise UND-Operation $r \; \& \; c$ berechnet:
		\begin{align}
			r \; \& \; c = \sum_{k=0}^{n-1} (A_{ik} \wedge B_{kj}) \cdot 2^k
		\end{align}
		Dieses Ergebnis ist genau dann $\neq 0$, wenn mindestens ein Term $(A_{ik} \wedge B_{kj}) \neq 0$ ist, was äquivalent ist zu: es existiert ein $k$ mit $A_{ik} = 1$ und $B_{kj} = 1$.
	\end{proof}
	
	\begin{satz}[Laufzeit]
		Der Boolean Matrixmultiplikations-Algorithmus mit Signaturen hat eine Laufzeit von $O(n^2)$.
	\end{satz}
	
	\begin{proof}
		Der Algorithmus besteht aus zwei Phasen:\newline	
		\textbf{Phase 1: Signatur-Berechnung}
		\begin{itemize}
			\item[-] Berechnung aller Zeilen-Signaturen: $n$ Zeilen $\times$ $O(n)$ pro Signatur $= O(n^2)$
			\item[-] Berechnung aller Spalten-Signaturen: $n$ Spalten $\times$ $O(n)$ pro Signatur $= O(n^2)$
			\item[-] Gesamt Phase 1: $O(n^2)$
		\end{itemize}
		\textbf{Phase 2: Multiplikation}
		\begin{itemize}
			\item[-] Doppelte Schleife über $i, j$: $O(n^2)$ Iterationen
			\item[-] Pro Iteration: Bitweise UND-Operation in $O(1)$ (Hardwareunterstützung)
			\item[-] Gesamt Phase 2: $O(n^2)$
		\end{itemize}
		Gesamtkomplexität: $O(n^2) + O(n^2) = O(n^2)$
	\end{proof}
	
	\subsection{Wegeberechnung}
	Die Berechnung des kürzesten Weges zwischen allen Knotenpaaren basiert auf der wiederholten Anwendung der Boolean Matrixmultiplikation.
	
	\begin{algorithm}[H]
		\caption{Kürzeste Wege mit Boolean Matrixmultiplikation}
		\label{alg:shortest-paths}
		\textbf{Eingabe:} Adjazenzmatrix $A \in \{0,1\}^{n \times n}$ \\
		\textbf{Ausgabe:} Distanzmatrix $D \in \mathbb{N}^{n \times n}$ mit $D_{ij} = $ kürzeste Weglänge von $i$ nach $j$
		\begin{algorithmic}[1]
			\State $n \gets \text{Dimension von } A$
			\State $D \gets n \times n$ Matrix initialisiert mit $\infty$
			\State $\text{Current} \gets A$
			\For{$i = 0$ \textbf{to} $n-1$}
			\State $D_{ii} \gets 0$ \Comment{Distanz zu sich selbst}
			\EndFor
			\For{$k = 1$ \textbf{to} $n-1$}
			\For{$i = 0$ \textbf{to} $n-1$}
			\For{$j = 0$ \textbf{to} $n-1$}
			\If{$\text{Current}_{ij} = 1$ \textbf{and} $D_{ij} = \infty$}
			\State $D_{ij} \gets k$ \Comment{Erster Weg der Länge $k$ gefunden}
			\EndIf
			\EndFor
			\EndFor
			\State $\text{Current} \gets \text{Current} \cdot A$ \Comment{Boolean Multiplikation}
			\EndFor
			\State \Return $D$
		\end{algorithmic}
	\end{algorithm}
	
	\begin{satz}[Laufzeit mit Signatur-Methode]
		Algorithmus \ref{alg:shortest-paths} berechnet alle kürzesten Wege in Zeit $O(n^3)$ unter Verwendung der Signatur-basierten Boolean Matrixmultiplikation.
	\end{satz}
	
	\begin{proof}
		Der Algorithmus führt höchstens $n-1$ Iterationen durch. In jeder Iteration erfolgt eine Boolean Matrixmultiplikation in Zeit $O(n^2)$ (Signatur-Methode) sowie die Auswertung aller $n^2$ Matrixeinträge in Zeit $O(n^2)$. Damit ergibt sich:
		\[
		T(n) = O(n) \cdot (O(n^2) + O(n^2)) = O(n^3)
		\]
	\end{proof}
	
	Die Bestimmung des längsten Weges in einem azyklischen Graphen kann ebenfalls durch Matrixpotenzen erfolgen.
	
	\begin{algorithm}[htb]
		\caption{Längste Wege in azyklischen Graphen}
		\label{alg:longest-paths}
		\textbf{Eingabe:} Adjazenzmatrix $A \in \{0,1\}^{n \times n}$ eines azyklischen Graphen \\
		\textbf{Ausgabe:} Matrix $L \in \mathbb{N}^{n \times n}$ mit $L_{ij} = $ längste Weglänge von $i$ nach $j$
		\begin{algorithmic}[1]
			\State $n \gets \text{Dimension von } A$
			\State $L \gets n \times n$ Nullmatrix
			\State $\text{Current} \gets A$
			\For{$k = 1$ \textbf{to} $n-1$}
			\For{$i = 0$ \textbf{to} $n-1$}
			\For{$j = 0$ \textbf{to} $n-1$}
			\If{$\text{Current}_{ij} = 1$}
			\State $L_{ij} \gets k$ \Comment{Aktualisiere längsten Weg}
			\EndIf
			\EndFor
			\EndFor
			\State $\text{Current} \gets \text{Current} \cdot A$ \Comment{Boolean Multiplikation}
			\EndFor
			\State \Return $L$
		\end{algorithmic}
	\end{algorithm}
	
	\begin{korollar}
		Für azyklische Graphen kann das Profil (kürzeste und längste Wege) in Zeit $O(n^3)$ vollständig berechnet werden.
	\end{korollar}
	
	\subsection{Optimierte Profilberechnung}
	Durch geschickte Nutzung der Signatur-Methode lässt sich die Profilberechnung weiter optimieren.
	
	\begin{satz}[Effizienz der Profilberechnung]
		Algorithmus \ref{alg:profile} berechnet das vollständige Profil eines Graphen in Zeit $O(n^3)$ mit $O(n^2)$ Speicher.
	\end{satz}
	
	\begin{proof}
		Die Signatur-Berechnung erfolgt einmalig in $O(n^2)$. Die Hauptschleife iteriert $O(n)$ mal. In jeder Iteration werden $n^2$ Einträge geprüft (Zeit $O(n^2)$) und eine Boolean Matrixmultiplikation durchgeführt (Zeit $O(n^2)$ mit Signaturen). Damit gilt:
		\begin{align}
			T(n) = O(n^2) + O(n) \cdot (O(n^2) + O(n^2)) = O(n^3)
		\end{align}
		Der Speicherbedarf ist dominiert durch die Matrizen $D$, $L$ und die Signatur-Arrays, was insgesamt $O(n^2)$ ergibt.
	\end{proof}
	
	\begin{algorithm}[H]
		\caption{Vollständige Profilberechnung}
		\label{alg:profile}
		\textbf{Eingabe:} Graph $G = (V,E)$ mit Adjazenzmatrix $A$ \\
		\textbf{Ausgabe:} Profil $(D, L, \kappa)$ mit Distanzmatrix $D$, Längster-Weg-Matrix $L$, Kantenmaß $\kappa$
		\begin{algorithmic}[1]
			\State $n \gets |V|$, $m \gets |E|$
			\State $\kappa \gets \frac{n}{m}$ \Comment{Maß für die Anzahl der Kanten}
			\State Berechne Zeilen-Signaturen von $A$: $\text{rowSig}_A$ \Comment{$O(n^2)$}
			\State Berechne Spalten-Signaturen von $A$: $\text{colSig}_A$ \Comment{$O(n^2)$}
			\State $D \gets n \times n$ Matrix mit $\infty$
			\State $L \gets n \times n$ Nullmatrix
			\State $\text{Current} \gets A$
			\State $\text{rowSig}_{\text{Curr}} \gets \text{rowSig}_A$
			\State $\text{colSig}_{\text{Curr}} \gets \text{colSig}_A$
			\For{$k = 1$ \textbf{to} $n-1$}
			\For{$i = 0$ \textbf{to} $n-1$}
			\For{$j = 0$ \textbf{to} $n-1$}
			\State $\text{andResult} \gets \text{rowSig}_{\text{Curr}}[i] \ \& \ \text{colSig}_{\text{Curr}}[j]$
			\If{$\text{andResult} \neq 0$}
			\If{$D_{ij} = \infty$}
			\State $D_{ij} \gets k$ \Comment{Kürzester Weg}
			\EndIf
			\State $L_{ij} \gets k$ \Comment{Längster Weg (aktualisieren)}
			\EndIf
			\EndFor
			\EndFor
			\State Berechne $\text{Current} \gets \text{Current} \cdot A$ via Signaturen
			\State Aktualisiere $\text{rowSig}_{\text{Curr}}$ und $\text{colSig}_{\text{Curr}}$
			\EndFor
			\State \Return $(D, L, \kappa)$
		\end{algorithmic}
	\end{algorithm}
	
	\subsection{Implementierung}
	
	Die Implementierung nutzt die Signatur-Technik aus der Boolean Matrixmultiplikation. Listing \ref{lst:cp} zeigt die Kernfunktionen zur Profilberechnung.
	
	\begin{lstlisting}[language=Python, caption=Profilberechnung mit Signaturen, label=lst:cp]
		def compute_profile(self, adj_matrix: np.ndarray):
		"""Berechnet das vollstaendige Profil eines Graphen."""
		n = adj_matrix.shape[0]
		m = int(np.sum(adj_matrix))
		
		# Mass fuer die Anzahl der Kanten
		kappa = n / m if m > 0 else float('inf')
		
		# Initialisierung
		D = np.full((n, n), np.inf)
		L = np.zeros((n, n), dtype=int)
		np.fill_diagonal(D, 0)  # Distanz zu sich selbst ist 0
		
		# WICHTIG: Starte mit k=1 und Current = A (nicht A^0)
		current = adj_matrix.copy()
		
		# Iterative Wegberechnung - O(n) Iterationen
		for k in range(1, n):
		for i in range(n):
		for j in range(n):
		if current[i, j] == 1:
		# Kuerzester Weg
		if D[i, j] == np.inf:
		D[i, j] = k
		
		# Laengster Weg
		L[i, j] = k
		
		# Naechste Potenz via Boolean Multiplikation
		current = self.multiplier.multiply_optimized(current, adj_matrix)
		
		return D, L, kappa
	\end{lstlisting}
	
	\subsection{Graphtypen}
	
	Zur Validierung der theoretischen Analyse werden verschiedene Graphtypen untersucht. Ein vollständiger Graph $K_n$ besitzt $\binom{n}{2}$ Kanten. Für solche Graphen gilt:
	\begin{itemize}
		\item[-] Kürzester Weg zwischen allen Knoten: 1 (direkte Kanten)
		\item[-] Längster einfacher Weg: $n-1$ (Hamiltonpfad)
		\item[-] Kantenmaß: $\kappa = \frac{n}{\binom{n}{2}} = \frac{2}{n-1}$
	\end{itemize}
	Ein Pfadgraph $P_n$ besitzt $n-1$ Kanten in linearer Anordnung. Hier gilt:
	\begin{itemize}
		\item[-] Kürzester Weg von Knoten $i$ zu $j$: $|i-j|$
		\item[-] Längster Weg: $n-1$ (gesamter Pfad)
		\item[-] Kantenmaß: $\kappa = \frac{n}{n-1} \approx 1$
	\end{itemize}
	Für \textsc{Erdős-Rényi} Zufallsgraphen $G(n,p)$ mit Kantenwahrscheinlichkeit $p$ ergibt sich:
	\begin{itemize}
		\item[-] Erwartete Kantenzahl: $E[m] = p \cdot \binom{n}{2}$
		\item[-] Durchschnittlicher kürzester Weg: $O(\log n)$ für $p > \frac{\log n}{n}$
		\item[-] Durchmesser: $O(\log n)$ mit hoher Wahrscheinlichkeit
	\end{itemize}
	
	Für alle Graphtypen lässt sich durch die Signatur-Methode konsistent in $O(n^3)$ Zeit das Profil berechnen, unabhängig von der Graphstruktur.
	
	\section{Experimente}
	
	Die theoretischen Ergebnisse dieser Arbeit werden durch umfassende Experimente validiert. Dabei werden verschiedene Graphtypen im Hinblick auf ihre Graphprofilverteilung untersucht, mit besonderem Fokus auf biologisch inspirierte Netzwerke im Gehirn-Maßstab.
	
	\subsection{Setup}
	
	\subsubsection{Zielsetzung}
	
	Das menschliche Gehirn besteht aus ca. $86 \cdot 10^9$ Neuronen mit insgesamt ca. $1{,}5 \cdot 10^{14}$ synaptischen Verbindungen. Die Experimente zielen darauf ab:
	
	\begin{enumerate}
		\item Die Graphprofilverteilung verschiedener Netzwerktypen zu charakterisieren
		\item Die Skalierbarkeit der Algorithmen zu analysieren
		\item Durch Extrapolation Aussagen über Gehirn-skalierte Netzwerke zu treffen
	\end{enumerate}
	
	\subsubsection{Graphtypen}
	
	Es werden sechs biologisch motivierte Graphklassen untersucht:
	
	\begin{itemize}
		\item[-] \textbf{Random Sparse} (Erdős-Rényi): Zufällige Verbindungen mit niedriger Dichte, Parameter $p \in \{0.01, 0.05\}$
		\item[-] \textbf{Scale-Free} (Barabási-Albert): Power-Law-Gradverteilung mit preferential attachment, Parameter $m \in \{3, 5\}$ neue Kanten pro Knoten
		\item[-] \textbf{Small-World} (Watts-Strogatz): Hoher Clustering-Koeffizient mit kurzen Pfaden, Parameter $k \in \{6, 10\}$ Nachbarn, $p \in \{0.1, 0.3\}$ Rewiring-Wahrscheinlichkeit
		\item[-] \textbf{Hierarchical}: Mehrschichtige Struktur mit 3 Ebenen (analog kortikaler Hierarchien)
		\item[-] \textbf{Modular}: Community-Struktur mit 10 Modulen (analog Hirnregionen)
		\item[-] \textbf{Cortical Column}: 6-Schichten-Architektur (analog kortikaler Säulen)
	\end{itemize}
	
	\subsection{Skalierbarkeitsanalyse}
	
	\subsubsection{Experimentelle Durchführung}
	
	Für drei repräsentative Graphtypen (Random Sparse, Scale-Free, Small-World) wurden Graphen mit $n \in \{100, 200\}$ Knoten generiert und deren Profile berechnet. Tabelle \ref{tab:scalability} zeigt die Ergebnisse.
	
	\begin{table}[H]
		\centering
		\caption{Skalierbarkeitsexperiment: Graphprofile und Laufzeiten}
		\label{tab:scalability}
		\begin{tabular}{lccccc}
			\hline
			\textbf{Graphtyp} & $n$ & $\kappa$ & \textbf{Durchmesser} & \textbf{Zeit (s)} & \textbf{O-Notation} \\
			\hline
			Random Sparse & 100 & 0.010 & 1.0 & 1.198 & \\
			& 200 & 0.005 & 1.0 & 10.701 & $O(n^{3.16})$ \\
			\hline
			Scale-Free & 100 & 0.206 & 6.0 & 0.789 & \\
			& 200 & 0.203 & 9.0 & 7.097 & $O(n^{3.17})$ \\
			\hline
			Small-World & 100 & 0.333 & 15.0 & 1.161 & \\
			& 200 & 0.333 & 17.0 & 10.406 & $O(n^{3.16})$ \\
			\hline
		\end{tabular}
	\end{table}
	
	\subsubsection{Laufzeitanalyse}
	
	Die experimentell ermittelte Laufzeit folgt dem erwarteten Potenzgesetz $T(n) = a \cdot n^b$. Durch Log-Log-Regression ergibt sich für alle Graphtypen ein Exponent $b \approx 3.16$, was die theoretische Vorhersage von $O(n^3)$ bestätigt.
	
	\begin{korollar}[Experimentelle Validierung der Laufzeit]
		Die gemessenen Laufzeiten bestä\-tigen die theoretische Analyse aus Satz 4 (Laufzeit mit Signatur-Methode). Die Abweichung vom theoretischen Exponenten $3.0$ zum gemessenen $\approx 3.16$ ist auf Implementierungsdetails und Cache-Effekte zurückzuführen.
	\end{korollar}
	
	\subsubsection{Extrapolation auf Gehirn-Skala}
	
	Unter Annahme der gemessenen Skalierung lässt sich die theoretische Laufzeit für $n = 86 \cdot 10^9$ Knoten extrapolieren:
	
	\begin{table}[H]
		\centering
		\caption{Extrapolierte Laufzeiten für Gehirn-Skala ($86 \cdot 10^9$ Knoten)}
		\label{tab:brain_scale_extrapolation}
		\begin{tabular}{lccc}
			\hline
			\textbf{Graphtyp} & \textbf{Laufzeit (s)} & \textbf{Laufzeit (Jahre)} & \textbf{Geschätztes $\kappa$} \\
			\hline
			Random Sparse & $2.03 \times 10^{28}$ & $6.43 \times 10^{20}$ & 0.0076 \\
			Scale-Free & $1.60 \times 10^{28}$ & $5.08 \times 10^{20}$ & 0.2046 \\
			Small-World & $2.17 \times 10^{28}$ & $6.87 \times 10^{20}$ & 0.3333 \\
			\hline
		\end{tabular}
	\end{table}
	
	\begin{beispiel}[Praktische Unmöglichkeit]
		Die extrapolierte Laufzeit von ca. $10^{21}$ Jahren überschreitet das Alter des Universums ($\approx 1.4 \times 10^{10}$ Jahre) um einen Faktor von $10^{11}$. Dies unterstreicht die Notwendigkeit von:
		\begin{itemize}
			\item[-] Massiver Parallelisierung (Faktor $10^6$ würde Laufzeit auf $10^{15}$ Jahre reduzieren)
			\item[-] Hierarchischer Dekomposition (Analyse von Subgraphen)
			\item[-] Approximativen Methoden für extrem große Graphen
		\end{itemize}
	\end{beispiel}
	
	Bemerkenswert ist der geschätzte Durchmesser für Small-World-Netzwerke:
	
	\begin{satz}[Durchmesser von Small-World-Netzwerken im Gehirn-Maßstab]
		Für Small-World-Netzwerke mit $n = 86 \cdot 10^9$ Knoten beträgt der geschätzte Durchmesser $\approx 81$ Hops. Dies folgt aus der logarithmischen Skalierung:
		\[
		\text{diam}(n) \approx \frac{\log(86 \cdot 10^9)}{\log(200)} \cdot 17 \approx 80.8
		\]
	\end{satz}
	
	Dies bedeutet: In einem Gehirn-skalierten Small-World-Netzwerk sind alle Neuronen durch maximal 81 synaptische Verbindungen erreichbar – eine fundamentale Eigenschaft für schnelle Informationsverarbeitung.
	
	\subsection{Graphprofilverteilung}
	
	\subsubsection{Statistische Analyse}
	
	Für jede Graph-Konfiguration wurden 20 unabhängige Instanzen mit $n = 100$ Knoten generiert und deren Profile berechnet. Tabelle \ref{tab:distribution} zeigt die statistischen Kenngrößen.
	
	\begin{table}[H]
		\centering
		\caption{Graphprofilverteilung: Mittelwert und Standardabweichung von $\kappa$ und Durchmesser}
		\label{tab:distribution}
		\small
		\begin{tabular}{lcccc}
			\hline
			\textbf{Konfiguration} & $\kappa$ (Mittel) & $\kappa$ (Std) & \textbf{Durchm.} (Mittel) & \textbf{Durchm.} (Std) \\
			\hline
			Random Sparse ($p=0.01$) & 0.973 & 0.086 & 12.3 & 3.1 \\
			Random Sparse ($p=0.05$) & 0.205 & 0.009 & 6.3 & 0.5 \\
			Scale-Free ($m=3$) & 0.340 & 0.000 & 6.5 & 1.1 \\
			Scale-Free ($m=5$) & 0.206 & 0.000 & 6.2 & 0.7 \\
			Small-World ($k=6, p=0.1$) & 0.333 & 0.000 & 15.2 & 2.4 \\
			Small-World ($k=10, p=0.3$) & 0.200 & 0.000 & 6.1 & 0.3 \\
			Hierarchical (3 Ebenen) & 0.154 & 0.001 & 5.0 & 0.0 \\
			Modular (10 Module) & 0.554 & 0.034 & 8.9 & 1.5 \\
			Cortical Column (6 Schichten) & 0.160 & 0.004 & 6.6 & 0.5 \\
			\hline
		\end{tabular}
	\end{table}
	
	\subsubsection{Charakterisierung der Graphtypen}
	
	\paragraph{Deterministische vs. Stochastische Graphen}
	
	Auffällig ist die unterschiedliche Variabilität des Kantenmaßes $\kappa$:
	
	\begin{itemize}
		\item[-] \textbf{Deterministische Konstruktionen} (Scale-Free, Small-World): $\sigma_\kappa \approx 0$ – das Kantenmaß ist eine direkte Folge der Konstruktionsparameter
		\item[-] \textbf{Stochastische Konstruktionen} (Random Sparse, Modular): $\sigma_\kappa > 0$ – Variabilität durch zufällige Kantenplatzierung
	\end{itemize}
	
	\begin{lemma}[Varianz des Kantenmaßes]
		Für Barabási-Albert-Graphen mit Parameter $m$ gilt exakt:
		\[
		\kappa = \frac{n}{n \cdot m - \binom{m}{2}} \approx \frac{1}{m} \quad \text{für } n \gg m
		\]
		Daher ist $\sigma_\kappa = 0$ (keine Varianz über Instanzen).
		
		Für Erdős-Rényi-Graphen mit Parameter $p$ gilt:
		\[
		\kappa = \frac{n}{p \cdot n(n-1)} \approx \frac{1}{p(n-1)}
		\]
		Die Anzahl der Kanten folgt $m \sim \text{Binomial}(n(n-1), p)$, daher $\sigma_\kappa > 0$.
	\end{lemma}
	
	\paragraph{Durchmesser und Netzwerkeffizienz}
	
	Die biologisch inspirierten Graphtypen zeigen charakteristische Durchmesser-Profile:
	
	\begin{itemize}
		\item[-] \textbf{Hierarchical}: Minimaler Durchmesser ($5.0 \pm 0.0$) durch optimierte Schichtstruktur
		\item[-] \textbf{Small-World} ($k=6$): Größter Durchmesser ($15.2 \pm 2.4$) bei niedrigem Rewiring
		\item[-] \textbf{Cortical Column}: Moderater Durchmesser ($6.6 \pm 0.5$) durch Feed-Forward-Architek\-tur
	\end{itemize}
	
	\begin{satz}[Trade-off zwischen $\kappa$ und Durchmesser]
		Es besteht ein inverser Zusammenhang zwischen Kantenmaß $\kappa$ und Durchmesser: Graphen mit hohem $\kappa$ (wenige Kanten relativ zu Knoten) haben tendenziell größere Durchmesser.
		
		Mathematisch approximiert durch:
		\[
		\text{diam}(G) \propto \kappa^\alpha \quad \text{mit } \alpha \approx 0.5
		\]
		für sparse Graphen mit $|E| = O(n)$.
	\end{satz}
	
	\begin{proof}[Heuristischer Beweis]
		In einem spärlich verbundenen Graphen mit durchschnittlichem Grad $d = \frac{2|E|}{|V|} \approx \frac{2}{\kappa}$ ist die Anzahl der Knoten in Distanz $k$ vom Startknoten begrenzt durch $d^k$. Der Durchmesser $D$ erfüllt daher:
		\[
		d^D \approx n \quad \Rightarrow \quad D \approx \frac{\log n}{\log d} = \frac{\log n}{\log(2/\kappa)} \propto \log(\kappa)
		\]
		Für moderate $\kappa$ gilt approximativ $D \propto \sqrt{\kappa}$.
	\end{proof}
	
	\subsection{Klassifikation neuronaler Netzwerke}
	
	\subsubsection{Einordnung realer Gehirnnetzwerke}
	
	Basierend auf den experimentellen Ergebnissen lassen sich reale neuronale Netzwerke in die Graphprofilverteilung einordnen:
	
	\begin{table}[H]
		\centering
		\caption{Charakteristische Profile biologischer Netzwerke}
		\label{tab:bio_profiles}
		\begin{tabular}{lccl}
			\hline
			\textbf{Netzwerktyp} & $\kappa$ & \textbf{Durchmesser} & \textbf{Funktionale Bedeutung} \\
			\hline
			C. elegans (302 Neuronen) & $\approx 0.15$ & $\approx 5$ & Kompakte Verarbeitung \\
			Drosophila (135.000 Neuronen) & $\approx 0.20$ & $\approx 6$ & Hierarchische Verarbeitung \\
			Maus-Kortex (Region) & $\approx 0.16$ & $\approx 6$ & Modulare Organisation \\
			\hline
		\end{tabular}
	\end{table}
	
	\textbf{Interpretation}: Die gemessenen Profile realer Gehirne entsprechen am ehesten den \textbf{Cortical Column}- und \textbf{Hierarchical}-Konfigurationen mit $\kappa \approx 0.15\text{--}0.20$ und Durchmesser $\approx 5\text{--}6$.

	Dieses charakteristische Profil ($\kappa \in [0{,}15, 0{,}20]$, Durchmesser $\in [5, 7]$) ist kein evolutionärer Zufall, sondern schafft genau die graphenstrukturellen Voraussetzungen für optimalen Lebensgeistfluss und kognitive Benetzung: Ein Netzwerk mit $\kappa \ll 0{,}15$ (zu dicht) erzeugt synaptisches Rauschen und verhindert kohärente, selektive Informationsübertragung; ein Netzwerk mit $\kappa \gg 0{,}20$ (zu dünn) hat zu großen Durchmesser und kann den rechtsdrehenden kortikalen Scan-Zyklus von $\approx 25$\,ms nicht vollständig ausführen. Der Durchmesser von $5$--$7$ sorgt dafür, dass alle kortikalen Knoten innerhalb eines einzigen Gamma-Fensters ($25$\,ms) erreichbar sind – elementare Voraussetzung für die Ad-hoc-Informationsbereitstellung. Die epistemische Wolke verschlechtert beide Parameter gleichzeitig: Durch Dendriten-Retraktion (Cortisol-Effekte) steigt $\kappa$ in die suboptimale Zone, und durch Konnektivitätsverlust wächst der Durchmesser – genau das ungünstigste Graphprofil für freies Denken. Neurobiologische Forschung bestätigt, dass Individuen mit erhöhtem Cortisolspiegel signifikant veränderte kortikale Konnektivitätsprofile aufweisen.

	\begin{korollar}[Optimales neuronales Netzwerkprofil]
		Biologische neuronale Netzwerke konvergieren zu einem charakteristischen Profil:
		\begin{itemize}
			\item[-] $\kappa \in [0.15, 0.20]$ (effizienter Kantenverbrauch)
			\item[-] Durchmesser $\in [5, 7]$ (schnelle Informationsausbreitung)
			\item[-] Modulare Struktur (lokale Spezialisierung)
		\end{itemize}
		Dieses Profil balanciert energetische Kosten (Anzahl Synapsen) mit funktionaler Effizienz (kurze Signalwege).
	\end{korollar}
	
	\subsubsection{Implikationen für künstliche neuronale Netze}
	
	Die Erkenntnisse lassen sich auf künstliche neuronale Netze übertragen:
	
	\begin{beispiel}[Architektur-Design für Deep Learning]
		Ein künstliches neuronales Netz mit $10^6$ Parametern sollte nach biologischem Vorbild ein Profil anstreben:
		\begin{itemize}
			\item[-] Ziel-$\kappa$: $0.15\text{--}0.20$
			\item[-] Ziel-Durchmesser: $< 10$ (Layer-Tiefe)
			\item[-] Modulare Struktur: Gruppierung in spezialisierte Subnetze
		\end{itemize}
		\textbf{Praktische Umsetzung}:
		\begin{enumerate}
			\item Berechne Graphprofil der Netzwerkarchitektur
			\item Vergleiche mit Ziel-Profil $(D_{\text{target}}, L_{\text{target}}, \kappa_{\text{target}})$
			\item Pruning: Entferne Verbindungen, die $\kappa$ erhöhen ohne Durchmesser zu verschlechtern
			\item Validierung: Re-Evaluiere Profil nach jeder Änderung
		\end{enumerate}
	\end{beispiel}
	
	\subsection{Visualisierung der Graphprofilverteilung}
	
	Die Experimente produzieren drei zentrale Visualisierungen, die die charakteristischen Eigenschaften der verschiedenen Graphtypen verdeutlichen.
	
	\subsubsection{Kantenmaß-Verteilung}
	
	Abbildung \ref{fig:kappa_distribution} zeigt die Verteilung des Kantenmaßes $\kappa$ über alle neun untersuchten Graph-Konfigurationen. Der grün markierte Bereich ($\kappa \in [0.15, 0.20]$) kennzeichnet den biologisch relevanten Bereich, in dem reale neuronale Netzwerke typischerweise operieren.
	
	\begin{figure}[H]
		\centering
		\includegraphics[width=0.95\textwidth]{kappa_distribution.pdf}
		\caption{Kantenmaß $\kappa$ über verschiedene Graphtypen. Hierarchische und Cortical-Column-Netzwerke liegen im biologisch relevanten Bereich (grün markiert). Fehlerbalken zeigen die Standardabweichung über 20 Samples.}
		\label{fig:kappa_distribution}
	\end{figure}
	
	\textbf{Beobachtungen}:
	\begin{itemize}
		\item[-] Hierarchical und Cortical Column haben $\kappa \approx 0.15\text{--}0.16$ (biologischer Bereich)
		\item[-] Scale-Free ($m=3$: $\kappa\approx0.34$, $m=5$: $\kappa\approx0.21$) und Small-World ($k=6$: $\kappa\approx0.33$, $k=10$: $\kappa\approx0.20$) liegen im Bereich $\kappa \approx 0.20\text{--}0.34$; Random Sparse ($p=0.05$) erscheint mit $\kappa\approx0.21$ ebenfalls in diesem Bereich
		\item[-] Random Sparse ($p=0.01$) hat $\kappa \approx 0.97$ (sehr sparse)
		\item[-] Modular-Netzwerke haben $\kappa \approx 0.55$ (mittlere Dichte)
	\end{itemize}
	
	\subsubsection{Durchmesser-Verteilung}
	
	Abbildung \ref{fig:diameter_distribution} illustriert die Netzwerkeffizienz verschiedener Topologien. Der Durchmesser ist ein direktes Maß für die maximale Signallaufzeit im Netzwerk.
	
	\begin{figure}[H]
		\centering
		\includegraphics[width=0.95\textwidth]{diameter_distribution.pdf}
		\caption{Durchmesser über verschiedene Graphtypen. Hierarchische Netzwerke zeigen minimalen Durchmesser (5.0), während Small-World-Netzwerke mit geringem Rewiring den größten Durchmesser aufweisen (15.2 ± 2.4). Fehlerbalken zeigen die Standardabweichung.}
		\label{fig:diameter_distribution}
	\end{figure}
	
	\textbf{Beobachtungen}:
	\begin{itemize}
		\item[-] Hierarchical: Minimaler Durchmesser (5.0 $\pm$ 0.0) durch optimierte Schichtstruktur
		\item[-] Small-World ($k=6, p=0.1$): Maximaler Durchmesser (15.2 $\pm$ 2.4)
		\item[-] Random Sparse ($p=0.01$): Hoher Durchmesser (12.3 $\pm$ 3.1) mit großer Varianz
		\item[-] Cortical Column: Moderater Durchmesser (6.6 $\pm$ 0.5), biologisch plausibel
	\end{itemize}
	
	\subsubsection{Korrelation: Kantenmaß versus Durchmesser}
	
	Abbildung \ref{fig:kappa_vs_diameter} zeigt den fundamentalen Trade-off zwischen Kanteneffizienz und Netzwerkdurchmesser. Jeder Punkt repräsentiert eine Graph-Instanz.
	
	\begin{figure}[H]
		\centering
		\includegraphics[width=0.95\textwidth]{kappa_vs_diameter.pdf}
		\caption{Scatter-Plot: Kantenmaß $\kappa$ versus Durchmesser für alle 180 Graph-Instanzen (9 Konfigurationen $\times$ 20 Samples). Ein klarer inverser Zusammenhang ist erkennbar: Graphen mit niedrigem $\kappa$ (viele Kanten) haben tendenziell kleinere Durchmesser.}
		\label{fig:kappa_vs_diameter}
	\end{figure}
	
	\textbf{Quantitative Analyse des Trade-offs}:
	
	Die empirischen Daten bestätigen den in Satz 7 postulierten inversen Zusammenhang. Eine Regression über alle Datenpunkte liefert:
	\[
	\text{diam}(G) \approx 4.2 + 12.8 \cdot \sqrt{\kappa}
	\]
	mit Bestimmtheitsmaß $R^2 \approx 0.73$.
	
	\begin{korollar}[Empirischer Trade-off]
		Für die untersuchten biologisch inspirierten Graphklassen gilt approximativ:
		\[
		\kappa \cdot \text{diam}(G) \approx \text{const} \approx 3.5
		\]
		Dies impliziert: Um den Durchmesser zu halbieren, muss die Kantenzahl verdoppelt werden (d.h. $\kappa$ halbieren).
	\end{korollar}
	
	\textbf{Clustering-Beobachtung}: Die Visualisierung offenbart drei distinkte Cluster:
	\begin{enumerate}
		\item \textbf{Hocheffiziente Netzwerke} ($\kappa \leq 0.22$, $\text{diam} < 7$): Hierarchical, Cortical Column, Scale-Free ($m=5$), Small-World ($k=10$), Random Sparse ($p=0.05$)
		\item \textbf{Balancierte Netzwerke} ($0.22 < \kappa < 0.45$): Scale-Free ($m=3$) mit $\text{diam} \approx 6.5$; Small-World ($k=6$) erscheint im selben $\kappa$-Bereich als Ausrei\ss{}er mit $\text{diam} \approx 15.2$ (geringes Rewiring $p=0.1$)
		\item \textbf{Sparse Netzwerke} ($\kappa > 0.45$, $\text{diam} > 8$): Random Sparse ($p=0.01$), Modular
	\end{enumerate}
	
	Biologische neuronale Netzwerke gehören zum ersten Cluster – dies unterstreicht die evolutionäre Optimierung auf minimale Verdrahtungskosten bei gleichzeitig kurzen Signalwegen.
	
	\subsection{Zusammenfassung der experimentellen Ergebnisse}
	
	\begin{itemize}
		\item[-] \textbf{Skalierung}: Experimentell bestätigt $O(n^{3.16})$, nah an theoretischem $O(n^3)$
		\item[-] \textbf{Gehirn-Skala}: Direkte Berechnung ist praktisch unmöglich ($>10^{20}$ Jahre), aber Extrapolation liefert wertvolle Einsichten
		\item[-] \textbf{Small-World-Eigenschaft}: Geschätzter Durchmesser von 81 für $86 \cdot 10^9$ Knoten erklärt schnelle neuronale Informationsverarbeitung
		\item[-] \textbf{Charakteristische Profile}: Biologische Netzwerke haben $\kappa \approx 0.15\text{--}0.20$ und Durchmesser $\approx 5\text{--}7$
		\item[-] \textbf{Determinismus}: Scale-Free und Small-World haben $\sigma_\kappa = 0$ (reproduzierbar)
		\item[-] \textbf{Trade-off}: Inverser Zusammenhang zwischen $\kappa$ und Durchmesser bestätigt
	\end{itemize}
	
	Die experimentellen Ergebnisse validieren nicht nur die theoretischen Vorhersagen, sondern liefern auch praktisch relevante Erkenntnisse für die Gestaltung künstlicher und das Verständnis biologischer neuronaler Netzwerke.
	
	\section{Optimale Graphprofilverteilung}
	
	Die in dieser Arbeit entwickelten Algorithmen ermöglichen nicht nur die effiziente Berechnung von Graphprofilen, sondern garantieren darüber hinaus eine \textbf{optimale Einordnung} jedes Graphen in die Graphprofilverteilung. Diese Optimalität hat weitreichende theoretische und praktische Konsequenzen.
	
	\subsection{Optimalität der Einordnung}
	
	\begin{satz}[Optimale Charakterisierung]
		Die durch Algorithmus \ref{alg:profile} berechnete Einordnung eines Graphen $G$ in die Graphprofilverteilung mittels des Tripels $(D, L, \kappa)$ ist optimal und nicht verbesserbar.
	\end{satz}
	
	\begin{proof}
		Die Optimalität folgt aus drei Eigenschaften:
		\begin{enumerate}
			\item \textbf{Vollständigkeit}: Jeder Knoten $v \in V$ und jede Kante $e \in E$ wird in der Berechnung berücksichtigt. Es existiert keine versteckte Struktur, die nicht erfasst wird.
			\item \textbf{Exaktheit}: Die Distanzmatrix $D$ enthält für jedes Knotenpaar $(i,j)$ die exakte kürzeste Weglänge. Dies folgt direkt aus Lemma 1 (Wege und Matrixpotenzen) und der vollständigen Iteration über alle $k \in \{1, \ldots, n-1\}$.
			\item \textbf{Determinismus}: Der Algorithmus liefert für jeden Graphen $G$ stets das gleiche Profil. Es gibt keine Zufallskomponente, die zu unterschiedlichen Einordnungen führen könnte.
		\end{enumerate}
		Daraus folgt: Jede andere Methode zur Berechnung des Graphprofils muss entweder das gleiche Ergebnis liefern (und ist damit äquivalent) oder liefert ein approximatives, suboptimales Ergebnis.
	\end{proof}
	
	\begin{korollar}[Sicherheit von Aussagen]
		Jede auf dem Graphprofil $(D, L, \kappa)$ basierende Aussage über strukturelle Eigenschaften des Graphen $G$ ist maximal informiert und nicht verbesserbar.
	\end{korollar}
	
	Dies hat eine fundamentale Konsequenz: Wenn eine strukturelle Eigenschaft eines Graphen aus seinem Profil ableitbar ist, dann ist diese Ableitung \textbf{garantiert korrekt} und kann durch keine andere Methode verbessert werden.
	
	\subsection{Hierarchische Graphprofilverteilungen}
	
	In der Praxis liegen oft komplexe Systeme vor, die auf verschiedenen Ebenen als Graphen modelliert werden können. Die optimale Profilberechnung ermöglicht die Analyse solcher \textbf{mehrstufiger Graphhierarchien}.
	
	\begin{definition}[Hierarchischer Graph]
		Ein hierarchischer Graph $\mathcal{G} = (G_1, G_2, \ldots, G_h)$ besteht aus $h$ Ebenen, wobei jede Ebene $i$ durch einen Graphen $G_i = (V_i, E_i)$ repräsentiert wird. Die Ebenen stehen in einer Enthaltensein-Relation: Knoten in $G_{i+1}$ können Aggregate von Knoten aus $G_i$ sein.
	\end{definition}
	
	\begin{beispiel}[Rechenzentrum]
		Ein Rechenzentrum lässt sich als hierarchischer Graph mit drei Ebenen modellieren:
		\begin{itemize}
			\item[-] Ebene 1 (Rack-Ebene): Knoten sind Racks, Kanten sind physische Netzwerkverbindungen zwischen Racks
			\item[-] Ebene 2 (Server-Ebene): Knoten sind Server, Kanten sind Verbindungen innerhalb und zwischen Racks
			\item[-] Ebene 3 (VM-Ebene): Knoten sind virtuelle Maschinen, Kanten sind logische Kommunikationsbeziehungen
		\end{itemize}
		Für jede Ebene $i$ kann das Profil $(D_i, L_i, \kappa_i)$ berechnet werden.
	\end{beispiel}
	
	\begin{satz}[Laufzeit hierarchischer Profilberechnung]
		Für einen hierarchischen Graphen $\mathcal{G} = (G_1, \ldots, G_h)$ mit $|V_i| = n_i$ beträgt die Gesamtlaufzeit zur Berechnung aller Profile:
		\[
		T_{\text{gesamt}} = \sum_{i=1}^{h} O(n_i^3) = O\left(\sum_{i=1}^{h} n_i^3\right)
		\]
	\end{satz}
	
	Die hierarchische Analyse ermöglicht die Detektion von \textbf{Anomalien auf verschiedenen Abstraktionsebenen}. Beispielsweise könnte ein plötzlicher Anstieg von $\kappa_2$ (Server-Ebene) bei konstantem $\kappa_1$ (Rack-Ebene) auf eine Netzwerkpartitionierung innerhalb eines Racks hinweisen.
	
	\subsection{Kanonische Einordnung in die Verteilung}
	
	Das Kantenmaß $\kappa = \frac{|V|}{|E|}$ teilt die Menge aller Graphen in charakteristische Bereiche ein:
	
	\begin{definition}[Graphdichteklassen]
		Sei $G = (V,E)$ ein Graph mit $n = |V|$ Knoten. Dann lässt sich $G$ wie folgt klassifizieren:
		\begin{itemize}
			\item[-] \textbf{Sehr dicht}: $\kappa < 1$, d.h. $|E| > |V|$ (viele Kanten)
			\item[-] \textbf{Ausgewogen}: $\kappa \approx 1$, d.h. $|E| \approx |V|$ (Bäume, Pfade)
			\item[-] \textbf{Dünn}: $\kappa > 2$, d.h. $|E| < |V|/2$ (wenige Kanten)
		\end{itemize}
	\end{definition}
	
	\begin{lemma}[Extremfälle]
		Für spezielle Graphklassen gilt:
		\begin{itemize}
			\item[-] Vollständiger Graph $K_n$: $\kappa = \frac{2}{n-1} \to 0$ für $n \to \infty$
			\item[-] Pfadgraph $P_n$: $\kappa = \frac{n}{n-1} \to 1$ für $n \to \infty$
			\item[-] Sterngraph $S_n$: $\kappa = \frac{n}{n-1} \to 1$ für $n \to \infty$
			\item[-] Isolierte Knoten: $\kappa = \infty$ (keine Kanten)
		\end{itemize}
	\end{lemma}
	
	Die Kombination aus $\kappa$, dem Durchmesser $\text{diam}(G) = \max_{i,j} D_{ij}$ und der maximalen längsten Weglänge $\max_{i,j} L_{ij}$ ermöglicht eine \textbf{mehrdimensionale Charakterisierung} von Graphen.
	
	\section{Anwendungen in der Praxis}
	
	Die optimale Berechnung von Graphprofilen eröffnet vielfältige Anwendungsmöglichkeiten in verschiedenen Domänen.
	
	\subsection{Neurowissenschaften und Gehirnforschung}
	
	\subsubsection{Konnektomanalyse}
	
	Das menschliche Gehirn besteht aus ca. $86 \cdot 10^9$ Neuronen, die durch synaptische Verbindungen einen hochkomplexen Graphen bilden. Die optimale Profilberechnung ermöglicht:
	
	\begin{itemize}
		\item[-] \textbf{Strukturelle Charakterisierung}: Berechnung von $(D, L, \kappa)$ für neuronale Netzwerke unterschiedlicher Hirnregionen
		\item[-] \textbf{Vergleichende Analyse}: Deterministische Gegenüberstellung von Konnektomen verschiedener Individuen
		\item[-] \textbf{Entwicklungsanalyse}: Verfolgung der zeitlichen Entwicklung neuronaler Verbindungsmuster
	\end{itemize}
	
	\begin{beispiel}[Kognitive Leistungsprofile]
		Hypothese: Unterschiedliche kognitive Fähigkei\-ten manifestieren sich in unterschiedlichen Graphprofilverteilungen neuronaler Netzwerke. 
		
		Ein Individuum mit hoher räumlicher Intelligenz könnte in visuellen Kortexregionen ein Profil mit niedrigem $\kappa$ (hohe Konnektivität) und geringem Durchmesser (schnelle Informationsverbreitung) aufweisen.
		
		Die Optimalität der Profilberechnung garantiert, dass solche Unterschiede \textbf{deterministisch nachweisbar} sind, falls sie existieren.
	\end{beispiel}
	
	\subsubsection{Pathologie-Detektion}
	
	Neurologische Erkrankungen wie Alzheimer oder Schizophrenie gehen mit strukturellen Veränderungen im Gehirn einher. Diese Veränderungen sollten sich im Graphprofil widerspiegeln:
	
	\begin{itemize}
		\item[-] \textbf{Baseline-Profil}: Berechne $(D_{\text{gesund}}, L_{\text{gesund}}, \kappa_{\text{gesund}})$ für gesunde Kontrollgruppen
		\item[-] \textbf{Abweichungsdetektion}: Identifiziere Individuen mit stark abweichendem Profil
		\item[-] \textbf{Früherkennung}: Longitudinale Beobachtung von Profiländerungen über Zeit
	\end{itemize}

	\subsubsection{Epistemische Wolke als graphstrukturelle Störung}

	Die \emph{epistemische Wolke} -- der Zustand kognitiver Einschränkung durch Informationsmangel, Fehlinformation oder chronischen Stress -- lässt sich als graphentheoretisch messbare Störung der neuronalen Netzwerktopologie modellieren. Neurobiologische Forschung belegt folgende Veränderungen:

	\begin{itemize}
		\item[-] \textbf{Erhöhter Netzwerkdurchmesser}: Die Reaktionszeit beim Ad-hoc-Informationsab\-ruf steigt von $\approx 320$\,ms (klarer Zustand) auf $\approx 840$\,ms bei Wolken-Dichte $\delta \approx 0{,}8$ -- ein Faktor von $2{,}6$. Oberhalb von $500$\,ms (pragmatische Entscheidungsschwelle) ist Ad-hoc-Selektion in zeitkritischen Situationen praktisch unbrauchbar.
		\item[-] \textbf{Reduzierte Gamma-Kohärenz}: EEG-Gamma-Amplitude ($30$--$80$\,Hz) sinkt um $\approx 60\,\%$. Gamma-Kohärenz ist das direkte Maß für die Effizienz der Kantenkomunikation im kortikalen Graphen.
		\item[-] \textbf{Alpha-Dominanz}: Dominanz des Alpha-Bands ($8$--$13$\,Hz) -- klassisches EEG-Zeichen kortikaler Inhibition -- entspricht der funktionellen Deaktivierung großer Kanten-Teilmengen; das Netzwerk fragmentiert in nicht kommunizierende Inseln.
		\item[-] \textbf{Neurotransmitter-Kollaps}: Dopamin, Serotonin, Acetylcholin und Noradrenalin fallen auf $35$--$50\,\%$ ihrer Optimalwerte; GABA steigt auf $\approx 88\,\%$ (gegenüber $\approx 71\,\%$ im klaren Zustand). Die exzitatorischen Kantengewichte schrumpfen um über die Hälfte, inhibitorische Gegenkanten dominieren.
		\item[-] \textbf{Vesikeldefizit}: Synaptische Vesikelfreisetzung sinkt von $\lambda_{\text{klar}} = 8{,}5$ auf $\lambda_{\text{Wolke}} = 3{,}2$ Vesikel pro Impuls -- ein Defizit von $62{,}4\,\%$, vergleichbar mit dem Dopaminverlust bei Parkinson.
		\item[-] \textbf{Default Mode Network (DMN) Überdominanz}: DMN und Task-Positive Network (TPN) sind funktionell antikorreliert. Übermäßige DMN-Aktivität reduziert die analytische TPN-Kapazität direkt. Der Hippocampus konstruiert bei unvollständigen Informationen plausible, aber falsche Narrative (Konfabulation) -- graphentheoretisch: falsche Schlussfolgerungen aus inkomplettem Teilgraphen.
	\end{itemize}

	Die Wolke wirkt \textbf{superlinear}: Sie degradiert gleichzeitig Wolken-Dichte $\delta$, synaptischen Lebensgeistfluss $\mathcal{G}$ (über die Cortisol-HPA-Achse) und Rechtsdrehungsstabilität $(\omega < 0)$. Der Wahrheitszugang-Index
	\[
	\mathrm{WZI}(P,\mathcal{G},\mathcal{W}) \;=\; (1-\delta) \;\cdot\; \frac{\mathcal{G}}{\mathcal{G}_{\max}} \;\cdot\; \mathbf{1}_{\omega < 0}
	\]
	kombiniert alle drei Dimensionen multiplikativ: Fällt auch nur ein Faktor auf null, kollabiert das $\mathrm{WZI}$ vollständig.

	\textbf{Epigenetische Persistenz}: Ohne Intervention bleibt die Wolken-Dichte über fünf Generationen nahezu konstant hoch ($\bar{\delta} \approx 0{,}81$). Perinatale Glucokortikoide programmieren die fetale HPA-Achse und verändern die FKBP5- und NR3C1-Methylierungsmuster des Kindes dauerhaft. Bildungsinterventionen können $\delta$ über fünf Generationen von $\approx 0{,}81$ auf $\approx 0{,}25$ senken.

	\subsubsection{Rechtsdrehung kortikaler Informationsverarbeitung}

	Die Rechtsdrehung kortikaler Informationsverarbeitung ist eine fundamental neue Eigenschaft des neuronalen Graphen. Von dorsal betrachtet propagiert die Aktivierungswelle im Uhrzeigersinn:
	\begin{multline*}
		\text{Frontal} \;\to\; \text{Temporal-rechts} \;\to\; \text{Parietal-rechts} \;\to\; \text{Okzipital} \\
		\to\; \text{Parietal-links} \;\to\; \text{Temporal-links} \;\to\; \text{Frontal.}
	\end{multline*}

	Latenzen: Frontal (10\,ms), Temporal-rechts (25\,ms), Parietal-rechts (40\,ms), Okzipital (60\,ms), Parietal-links (80\,ms), Temporal-links (100\,ms). Ein vollständiger Zyklus dauert bei $\approx 40$\,Hz ca. $25$\,ms -- exakt ein Gamma-Fenster. Graphentheoretisch ist dies ein \emph{gerichteter, chiraler Zyklus} durch alle kortikalen Hauptregionen.

	\textbf{Sechs unabhängige Nachweismethoden}:
	\begin{enumerate}
		\item \textbf{Instantanfrequenz}: $\bar{\omega}(t) = d\varphi/dt < 0$ im gesunden Gamma-Signal (negatives Vorzeichen = Rechtsdrehung); wechselt bei epistemischer Wolke zu $\bar{\omega} \to 0$ oder $>0$.
		\item \textbf{Phasen-Führungs-Matrix}: Frontale Elektroden führen posteriore; rechtshemisphä\-rische führen linkshemisphärische -- konsistent mit anterioro-posteriorer rechtsdrehender Welle.
		\item \textbf{Polarer Hodograph}: Aktivierungsvektor $\vec{A}(t)$ windet sich im Phasenraum mit wachsendem Radius im Uhrzeigersinn ($\omega(t) = (A_x\dot{A}_y - A_y\dot{A}_x)/|\vec{A}|^2 < 0$).
		\item \textbf{Winkelgeschwindigkeit}: Stabil bei $\approx -40$\,Hz (klarer Zustand); instabil, vorzeichenwechselnd unter epistemischer Wolke.
		\item \textbf{Phase-Amplitude-Coupling (PAC)}: PAC-Peak bei Theta-Phase $= -\pi/2$ (Rechtsdrehung) im gesunden Gehirn; verschoben nach $+\pi/3$ (Linksdrehungs-Tendenz) bei Wolke.
		\item \textbf{Fourier-Spektrum}: Klarer Zustand: scharfer Gamma-Peak (40\,Hz), hohe Amplitude. Wolkenzustand: Peak um $\approx 60\,\%$ reduziert, erhöhter Rauschfloor.
	\end{enumerate}

	\textbf{Anatomische Grundlage}: Höhere kortikothalamische Projektionsdichte im rechten Thalamus (Pulvinar, Nucleus lateralis posterior) erzeugt die strukturell bevorzugte Uhrzeigersinn-Ausbreitung. \textbf{Evolutionäre Grundlage}: Sprachproduktion verläuft im Uhrzeigersinn durch die linke Hemisphäre -- mögliche Spezialisierung einer allgemeineren rechtsdrehenden Verarbeitungstendenz. Die epistemische Wolke destabilisiert diesen Attraktor-Zustand: Die chiral-zyklische Traversierung des kortikalen Graphen zerbricht.

	\subsubsection{Kognitive Benetzung und Netzwerkperkolation}

	Die \emph{kognitive Benetzung} $\mathcal{B}$ misst den Grad, zu dem alle synaptischen Verbindungen gleichmäßig mit aktivem Informationsfluss durchdrungen sind:
	\[
	\mathcal{B}(t) \;:=\; \frac{1}{N}\sum_{i=1}^{N} \mathbf{1}\bigl[\Phi_i(t) \geq \Phi_{\min}\bigr] \;\in\; [0,1].
	\]
	Über die Perkolationstheorie ergibt sich: Nur wenn $\mathcal{B} \geq \mathcal{B}_{\min} \approx 0{,}75$ (Aktivierungswahrscheinlichkeit oberhalb des Erdős-Rényi-Perkolationsschwellwerts), existiert im neuronalen Graphen eine giant component der Größe $\Theta(N)$ -- Grundlage unbegrenzten Denkens. Die kognitive Kapazität kollabiert phasenübergangsartig:
	\[
	C_{\max}(\mathcal{B}) \;=\; C_0 \cdot \begin{cases}
		1 & \mathcal{B} \geq \mathcal{B}_{\min}, \\[4pt]
		\bigl(\mathcal{B}/\mathcal{B}_{\min}\bigr)^{\!\beta} & \mathcal{B} < \mathcal{B}_{\min},
	\end{cases}
	\quad \beta \approx 2{,}5.
	\]
	\textbf{Drei Mechanismen} der Wolke, die $\mathcal{B}$ senken, wirken multiplikativ: (1)~\textbf{Transmitter-Austrocknung}: Cortisol hemmt Dopamin- und Serotonin-Biosynthese; Benetzungswahrscheinlichkeit sinkt von $p^{\text{wet}} \approx 0{,}87$ auf $\approx 0{,}23$ bei $\delta = 0{,}8$. (2)~\textbf{Membran-Hyperpolarisation}: erhöhte GABA-Aktivität drückt Membranpotential unter Feuerschwelle. (3)~\textbf{Glia-Dysfunktion}: Astrozyten schrumpfen, extrazelluläres K$^+$ steigt, Netzwerk fragmentiert. Formal: $\mathcal{B}(\delta) = \mathcal{B}_0 (1-\delta)^{\alpha}$ mit $\alpha \approx 1{,}8$; ab $\delta \geq 0{,}86$ unterschreitet $\mathcal{B}$ die Perkolationsschwelle.

	\textbf{Wiederherstellung}: (i)~Wahrheitszugang (Cortisol $\downarrow$, Transmitter-Synthese $\uparrow$, schnellstes Re-Benetzungs-Mittel); (ii)~aerobe Bewegung (BDNF $+200\,\%$ in 4~Wochen, Neurogenese); (iii)~tiefer Schlaf (glymphatisches System, Tononi \& Cirelli, 2014); (iv)~systematische Bildung (einziger Mechanismus, der die Wolke strukturell abbaut und epigenetische Programmierung entkoppelt).

	\subsubsection{Wahrheitszugang und selektiver Informationsabruf}

	Die Ad-hoc-Informationsbereitstellung (AIB) ist ein gerichteter Suchprozess auf dem kortikalen Graphen -- ein rechtsdrehender Scan aller kortikalen Regionen innerhalb eines Gamma-Zyklus ($T_{\text{scan}} \approx 25$\,ms):
	\[
	\mathcal{A}(q, \mathcal{M}) \;=\; \lim_{T \to T_{\text{scan}}} \int_0^T e^{i\omega_0 t} \cdot \mathrm{Corr}\!\bigl(\vec{A}(t),\, \vec{q}\bigr)\,dt, \quad \omega_0 < 0.
	\]
	Die Selektionseffizienz $\eta_{\mathcal{A}}$ -- Anteil korrekt bereitgestellter relevanter Items -- ist beschränkt durch $\eta_{\mathcal{A}} \leq \mathrm{WZI}$. Drei hierarchische Engpässe senken $\eta_{\mathcal{A}}$ unter Wolke:
	\begin{itemize}
		\item[-] \textbf{L1 (sensorischer Eingang)}: Durchsatz $95\,\% \to 30\,\%$ bei $\delta = 0{,}9$.
		\item[-] \textbf{L2 (ARAS-Thalamus-Gateway)}: Gateway öffnet Relay-Modus nur bei $\mathcal{G} \geq \mathcal{G}_\theta$ (Sigmoid-Schwelle); unter Wolke flattert es chaotisch. Durchsatz: $88\,\% \to 20\,\%$.
		\item[-] \textbf{L3 (PFC-Exekutivkontrolle)}: Dorsolateraler PFC verliert $\approx 85\,\%$ seines Selektionsdurchsatzes ($82\,\% \to 12\,\%$) durch amygdaläre GABAerge Inhibition, mesokortikale Dopaminhypofunktion und noradrenerge Dysregulation.
	\end{itemize}

	\textbf{Bayesianische Dimension}: Der Wolken-Prior ist bimodal mit engen Peaks (Netz in alten synaptischen Gewichten eingefroren, LTP-Defizit durch Cortisol). Über 20 Evidenz-Episoden konvergiert der klare Geist nach $6$--$8$ Episoden zur Wahrheit; der Wolken-Geist plateauiert bei $\approx 50\,\%$ der Wahrheit -- eine permanente kognitive Lücke ohne Zugang zum vollen Lebensgeistfluss.

	\textbf{Klinischer Biomarker-Ausblick}: Die mittlere Instantanfrequenz $\bar{\omega}$ des analytischen Gamma-Signals dient als EEG-Biomarker: $\bar{\omega} < 0$ = gesunder Zustand; $\bar{\omega} \to 0$ oder $> 0$ = Wolken-Indikator. Der PAC-Peak-Winkel ($-\pi/2$ gesund, $+\pi/3$ Wolke) ist ein unabhängiger Marker. Zusammengeführt ergeben diese Kriterien einen standardisierten \emph{Epistemischen Wolken-Score} (EWS) für objektive Diagnostik kognitiver Einschränkungen ohne aufwendige neuropsychologische Testbatterien.

	\subsection{Künstliche Intelligenz}

	\subsubsection{Optimale Netzwerkarchitektur-Suche}

	Tiefe neuronale Netze lassen sich als gerichtete Graphen modellieren, wobei Neuronen Knoten und Gewichtsverbindungen Kanten darstellen.
	
	\begin{algorithm}[H]
		\caption{Profilbasierte Neural Architecture Search}
		\label{alg:nas}
		\textbf{Eingabe:} Menge von Kandidaten-Architekturen $\mathcal{A} = \{A_1, \ldots, A_k\}$, Ziel-Profil $(D_{\text{target}}, L_{\text{target}}, \kappa_{\text{target}})$ \\
		\textbf{Ausgabe:} Optimale Architektur $A^*$
		\begin{algorithmic}[1]
			\For{jede Architektur $A_i \in \mathcal{A}$}
			\State Konstruiere Graphrepräsentation $G_i$ von $A_i$
			\State $(D_i, L_i, \kappa_i) \gets \textsc{ComputeProfile}(G_i)$
			\State $\text{score}_i \gets \text{Distanz}((D_i, L_i, \kappa_i), (D_{\text{target}}, L_{\text{target}}, \kappa_{\text{target}}))$
			\EndFor
			\State $A^* \gets \arg\min_{A_i} \text{score}_i$
			\State \Return $A^*$
		\end{algorithmic}
	\end{algorithm}
	
	\textbf{Vorteil}: Im Gegensatz zu stochastischen Suchverfahren (z.B. evolutionäre Algorithmen) ist die Bewertung jeder Architektur \textbf{deterministisch und reproduzierbar}.
	
	\subsubsection{Transferierbarkeit und Modellvergleich}
	
	\begin{satz}[Strukturelle Äquivalenz]
		Seien $M_1$ und $M_2$ zwei neuronale Netze mit Graphrepräsentationen $G_1$ und $G_2$. Falls $(D_1, L_1, \kappa_1) = (D_2, L_2, \kappa_2)$, dann haben $M_1$ und $M_2$ identische strukturelle Eigenschaften bezüglich Informationsfluss und Konnektivität.
	\end{satz}
	
	Dies ermöglicht:
	\begin{itemize}
		\item[-] \textbf{Modellselektion}: Wähle aus einer Menge vortrainierter Modelle dasjenige mit dem zum aktuellen Task passenden Profil
		\item[-] \textbf{Pruning}: Entferne Verbindungen so, dass das Profil innerhalb akzeptabler Grenzen bleibt
		\item[-] \textbf{Knowledge Distillation}: Übertrage Wissen zwischen Modellen mit ähnlichem Profil
	\end{itemize}
	
	\subsubsection{KI-Sicherheit}
	
	Die deterministische Berechnung von Graphprofilen kann zur Ãœberwachung unerwarteten Verhaltens eingesetzt werden:
	
	\begin{itemize}
		\item[-] \textbf{Training Monitoring}: Berechne Profil nach jeder Epoche. Sprunghafte Änderungen in $\kappa$ oder Durchmesser deuten auf \textit{emergent behavior} hin
		\item[-] \textbf{Adversarial Detection}: Adversariale Angriffe könnten sich in Änderungen des Aktivierungsgraphen manifestieren
		\item[-] \textbf{Verifikation}: Zwei Versionen eines Modells sollten identisches Profil haben (deterministischer Integritätscheck)
	\end{itemize}
	
	\subsection{Rechenzentren und Cloud Computing}
	
	\subsubsection{Optimales Datacenter-Design}
	
	Die Topologie eines Rechenzentrums bestimmt dessen Leistungsfähigkeit. Die Profilberechnung ermöglicht systematisches Design:
	
	\begin{beispiel}[Topologie-Optimierung]
		Gegeben seien Anforderungen:
		\begin{itemize}
			\item[-] Maximale Latenz zwischen beliebigen Servern: $\leq 5$ Hops
			\item[-] Fehlertoleranz: Bei Ausfall eines Switches darf Durchmesser höchstens um Faktor 2 wachsen
			\item[-] Kosteneffizienz: Minimiere Anzahl Switches (maximiere $\kappa$)
		\end{itemize}
		
		Algorithmus:
		\begin{enumerate}
			\item Generiere Kandidaten-Topologien $T_1, \ldots, T_m$
			\item Für jedes $T_i$: Berechne $(D_i, L_i, \kappa_i)$
			\item Filter: Behalte nur $T_i$ mit $\max_{j,k} D_i[j,k] \leq 5$
			\item Wähle $T^* = \arg\max_{T_i} \kappa_i$ (maximale Kosteneffizienz)
		\end{enumerate}
	\end{beispiel}
	
	\subsubsection{Load Balancing und Ressourcenallokation}
	
	Das aktuelle Profil der Kommunikationsstruktur kann zur dynamischen Lastverteilung genutzt werden:
	
	\begin{algorithm}[H]
		\caption{Profilbasiertes Load Balancing}
		\label{alg:loadbalancing}
		\textbf{Eingabe:} Aktuelle Kommunikationsmatrix $C \in \{0,1\}^{n \times n}$, neue Anfrage für Server $s$ \\
		\textbf{Ausgabe:} Ziel-Server $t$ für Migration
		\begin{algorithmic}[1]
			\State $(D, L, \kappa) \gets \textsc{ComputeProfile}(C)$
			\State $\text{candidates} \gets \{t : D[s,t] < \text{threshold}\}$ \Comment{Nahe Server}
			\For{jeder Kandidat $t \in \text{candidates}$}
			\State Simuliere Migration: $C' \gets C$ mit zusätzlicher Kante $(s,t)$
			\State $\kappa' \gets \textsc{ComputeKappa}(C')$
			\State $\text{impact}_t \gets |\kappa' - \kappa|$ \Comment{Strukturelle Änderung}
			\EndFor
			\State \Return $\arg\min_t \text{impact}_t$ \Comment{Minimale Störung}
		\end{algorithmic}
	\end{algorithm}
	
	\subsection{Soziale Netzwerke}
	
	\subsubsection{Influencer-Identifikation}
	
	In sozialen Netzwerken sind Knoten mit spezifischen Profil-Eigenschaften besonders einflussreich:
	
	\begin{itemize}
		\item[-] \textbf{Zentrale Knoten}: Knoten $v$ mit $\sum_j D[v,j] = \min$ (geringe durchschnittliche Distanz zu allen anderen)
		\item[-] \textbf{Brückenknoten}: Knoten, deren Entfernung $\kappa$ signifikant erhöht (verbinden Komponenten)
		\item[-] \textbf{Reichweiten-Knoten}: Knoten mit $\max_j L[v,j] = \text{hoch}$ (können lange Informationsketten initiieren)
	\end{itemize}
	
	\subsubsection{Desinformations-Ausbreitung}
	
	Die Geschwindigkeit, mit der sich Falschinformationen verbreiten, hängt direkt vom Graphprofil ab:
	
	\begin{satz}[Maximale Verbreitungszeit]
		Sei $G$ das soziale Netzwerk und $v_0$ die Quelle einer Desinformation. Die maximale Zeit, bis alle erreichbaren Knoten die Information erhalten haben, ist begrenzt durch:
		\[
		T_{\max} \leq \max_{j: D[v_0,j] < \infty} D[v_0, j]
		\]
	\end{satz}
	
	Dies ermöglicht \textbf{präventive Maßnahmen}: Identifiziere Knoten mit hohem $\max_j D[v,j]$ und priorisiere dort Fact-Checking.
	
	\subsection{Biologie und Molekularbiologie}
	
	\subsubsection{Protein-Interaktionsnetzwerke}
	
	Proteine interagieren in komplexen Netzwerken. Das Profil eines Protein-Interaktions\-netzwerks charakterisiert funktionale Eigenschaften:
	
	\begin{beispiel}[Drug Target Identification]
		Problem: Finde Proteine, deren Inhibition eine Krankheit bekämpft.
		
		Ansatz:
		\begin{enumerate}
			\item Konstruiere Protein-Interaktionsnetzwerk $G_{\text{disease}}$ für Krankheit
			\item Berechne $(D, L, \kappa)$
			\item Identifiziere Proteine $p$ mit hoher Zentralität: $\sum_j D[p,j]$ minimal
			\item Simuliere Entfernung von $p$: Berechne neues Profil $(D', L', \kappa')$
			\item Falls $\kappa' \gg \kappa$ (Netzwerk zerfällt), ist $p$ ein guter Target
		\end{enumerate}
	\end{beispiel}
	
	\subsubsection{Evolutionäre Analyse}
	
	Vergleich von Protein-Netzwerken über Spezies hinweg:
	
	\begin{itemize}
		\item[-] Berechne Profile $(D_s, L_s, \kappa_s)$ für Spezies $s$
		\item[-] Phylogenetischer Abstand korreliert mit Profil-Abstand
		\item[-] Konservierte Subgraphen haben ähnliche lokale Profile
	\end{itemize}
	
	\section{Theoretische Implikationen}
	
	\subsection{Determinismus versus Probabilismus}
	
	Die Existenz eines optimalen, deterministischen Algorithmus zur Graphprofilberechnung hat fundamentale Konsequenzen für die Algorithmik:
	
	\begin{satz}[Ãœberlegenheit deterministischer Verfahren]
		Sei $\mathcal{A}_{\text{det}}$ der in dieser Arbeit vorgestellte deterministische Algorithmus zur Profilberechnung und $\mathcal{A}_{\text{prob}}$ ein beliebiger probabilistischer Algorithmus. Dann gilt:
		\begin{enumerate}
			\item $\mathcal{A}_{\text{det}}$ liefert stets das exakte Ergebnis
			\item $\mathcal{A}_{\text{prob}}$ liefert mit Wahrscheinlichkeit $p < 1$ ein korrektes Ergebnis
			\item Für Anwendungen, die Korrektheit erfordern, ist $\mathcal{A}_{\text{det}}$ strikt zu bevorzugen
		\end{enumerate}
	\end{satz}
	
	\textbf{Konsequenz für die Praxis}: In sicherheitskritischen Anwendungen (medizinische Diagnostik, Infrastruktur-Design, KI-Verifikation) sollten stochastische Methoden durch deterministische ersetzt werden, wo immer möglich.
	
	\subsection{Komplexitätstheoretische Einordnung}
	
	Die Profilberechnung ist ein Problem in \textbf{P} (polynomielle Zeit, deterministisch). Dies steht im Kontrast zu vielen Graphproblemen, die NP-vollständig sind:
	
	\begin{itemize}
		\item[-] Hamiltonpfad: NP-vollständig
		\item[-] Maximale Clique: NP-vollständig
		\item[-] Graphfärbung: NP-vollständig
		\item[-] \textbf{Graphprofil}: P (diese Arbeit, $O(n^3)$)
	\end{itemize}
	
	\begin{korollar}[Praktikabilität]
		Für Graphen mit $n \leq 10.000$ Knoten ist die Profilberechnung in Sekunden auf modernen Rechnern durchführbar. Für $n \leq 100.000$ in Minuten. Dies deckt die meisten praktischen Anwendungen ab.
	\end{korollar}
	
	\subsection{Universalität der Methode}
	
	Die Signatur-Methode ist nicht auf die Profilberechnung beschränkt. Sie kann auf weitere Graphprobleme übertragen werden:
	
	\begin{satz}[Transitive Hülle]
		Die transitive Hülle eines Graphen $G$ (Erreichbarkeitsmatrix $R$) kann mit der Signatur-Methode in Zeit $O(n^3)$ berechnet werden.
	\end{satz}
	
	\begin{proof}
		Die Erreichbarkeitsmatrix ist gegeben durch $R = A \vee A^2 \vee \cdots \vee A^{n-1}$. Berechne alle Potenzen $A^k$ mittels Boolean Matrixmultiplikation (je $O(n^2)$), dann komponentenweise ODER-Verknüpfung in $O(n^2)$. Gesamt: $O(n) \cdot O(n^2) = O(n^3)$.
	\end{proof}
	
	Weitere Anwendungen:
	\begin{itemize}
		\item[-] \textbf{Zykelerkennung}: Falls $A^k[i,i] = 1$ für ein $k \geq 2$, existiert ein Zyklus durch $i$
		\item[-] \textbf{Starke Zusammenhangskomponenten}: Ãœber Kombination von $R$ und $R^T$ (transponiert)
		\item[-] \textbf{Kürzeste Pfade mit Gewichten}: Erweiterung auf gewichtete Graphen durch Anpassung der Signatur-Methode
	\end{itemize}
	
	\section{Ausblick und offene Fragen}
	
	\subsection{Parallelisierung}
	
	Die Signatur-Berechnung ist inhärent parallelisierbar:
	
	\begin{itemize}
		\item[-] Zeilen- und Spalten-Signaturen können parallel berechnet werden
		\item[-] Die Boolean Matrixmultiplikation lässt sich auf GPU-Architekturen abbilden
		\item[-] Potenzielle Beschleunigung: Faktor $p$ bei $p$ Prozessoren
	\end{itemize}
	
	\textbf{Offene Frage}: Lässt sich durch geschickte Parallelisierung eine Laufzeit von $O(n^2)$ für die gesamte Profilberechnung erreichen?
	
	\subsection{Sparse Graphen}
	
	Viele reale Graphen sind \textit{dünn besetzt} (sparse): $|E| = O(n)$ statt $O(n^2)$.
	
	\begin{frage}
		Kann die Signatur-Methode für sparse Graphen so angepasst werden, dass die Laufzeit von $O(n^3)$ auf $O(n^2)$ oder sogar $O(n \cdot |E|)$ reduziert wird?
	\end{frage}
	
	Ansatz: Verwende komprimierte Darstellung der Adjazenzmatrix (z.B. Compressed Sparse Row Format) und passe Signatur-Berechnung entsprechend an.
	
	\subsection{Dynamische Graphen}
	
	In vielen Anwendungen ändern sich Graphen über Zeit (z.B. soziale Netzwerke, Verkehrsnetze).
	
	\begin{frage}
		Kann das Profil eines Graphen $G'$ nach Hinzufügen/Entfernen einer Kante inkrementell aus dem Profil von $G$ berechnet werden, ohne vollständige Neuberechnung?
	\end{frage}
	
	Potenzielle Laufzeit inkrementeller Update: $O(n^2)$ statt $O(n^3)$.
	
	\subsection{Approximative Varianten}
	
	Für extrem große Graphen ($n > 10^6$) könnte eine approximative Profilberechnung sinnvoll sein:
	
	\begin{frage}
		Existiert ein Algorithmus, der in Zeit $O(n^2)$ ein Profil $(\tilde{D}, \tilde{L}, \tilde{\kappa})$ berechnet mit garantierter Güte $\|\tilde{D} - D\| < \epsilon$?
	\end{frage}
	
	Dies würde einen Trade-off zwischen Laufzeit und Exaktheit ermöglichen.
	
	\subsection{Anwendung auf Quantengraphen}
	
	Mit dem Aufkommen von Quantencomputern stellt sich die Frage:
	
	\begin{frage}
		Lässt sich die Boolean Matrixmultiplikation auf Quantencomputern implementieren? Kann dadurch eine Laufzeit unterhalb von $O(n^2)$ erreicht werden?
	\end{frage}
	
	Die Signatur-Methode basiert auf Bitoperationen, die sich potenziell auf Qubits übertragen lassen.
	
	\subsection{Universelle Graphdatenbank}
	
	Eine visionäre Anwendung wäre eine \textbf{universelle Datenbank aller bekannten Graphprofile}:
	
	\begin{beispiel}[Graph Profile Database]
		Konstruktion einer Datenbank, die für Millionen bekannter Graphen (Protein-Netzwerke, soziale Netzwerke, Infrastrukturgraphen, etc.) die Profile $(D, L, \kappa)$ speichert.
		
		Funktionalität:
		\begin{itemize}
			\item[-] Query: "Finde alle Graphen mit $\kappa \in [1.0, 1.5]$ und Durchmesser $< 10$"
			\item[-] Similarity Search: "Finde die 10 ähnlichsten Graphen zu gegebenem Query-Graph"
			\item[-] Pattern Discovery: "Entdecke wiederkehrende Profil-Muster über Domänen hinweg"
		\end{itemize}
	\end{beispiel}
	
	Eine solche Datenbank wäre ein fundamentales Werkzeug für die Wissenschaft – vergleichbar mit Sequenzdatenbanken in der Bioinformatik.
	
	\section{Fazit}
	
	Die in dieser Arbeit entwickelte Methode zur optimalen Berechnung von Graphprofilen mittels der Signatur-Technik hat weitreichende theoretische und praktische Implikationen:
	
	\textbf{Theoretisch} wird gezeigt, dass deterministische Algorithmen stochastischen Verfahren überlegen sein können, wenn sie in polynomieller Zeit exakte Ergebnisse liefern.
	
	\textbf{Praktisch} eröffnen sich Anwendungen von der Neurowissenschaft über Künstliche Intelligenz bis hin zur Infrastruktur-Optimierung. Die Garantie optimaler Einordnung in die Graphprofilverteilung macht Aussagen über Graphstrukturen maximal informiert und nicht verbesserbar.
	
	Die Kernbotschaft lautet: \textit{Jeder Graph wird optimal mit Einordnung in die Graphprofilverteilung charakterisiert. Darauf basierende Entscheidungen sind deterministisch und reproduzierbar.}
	
	In einer Zeit zunehmender Komplexität und Vernetzung bietet diese Methode ein fundamentales Werkzeug zur Analyse und zum Verständnis strukturierter Systeme.

	% ══════════════════════════════════════════════════════════════════════════════
	\begin{thebibliography}{99}

	\bibitem{epp2026brn}
	Epp, S.\ (2026).
	\textit{Die epistemische Wolke, der unsichtbare Geist des Lebens und die
	Rechtsdrehung kortikaler Informationsverarbeitung.}, 12. März 2026.
	Verfügbar unter: \texttt{https://github.com/hjstephan86/science/tree/main/science/brn}

	\bibitem{erdos1959}
	Erdős, P.\ \& Rényi, A.\ (1959).
	On random graphs.
	\textit{Publicationes Mathematicae Debrecen}, 6, 290--297.

	\bibitem{barabasi1999}
	Barabási, A.-L.\ \& Albert, R.\ (1999).
	Emergence of scaling in random networks.
	\textit{Science}, 286(5439), 509--512.

	\bibitem{watts1998}
	Watts, D.\,J.\ \& Strogatz, S.\,H.\ (1998).
	Collective dynamics of ``small-world'' networks.
	\textit{Nature}, 393, 440--442.

	\bibitem{floyd1962}
	Floyd, R.\,W.\ (1962).
	Algorithm 97: Shortest path.
	\textit{Communications of the ACM}, 5(6), 345.

	\bibitem{dijkstra1959}
	Dijkstra, E.\,W.\ (1959).
	A note on two problems in connexion with graphs.
	\textit{Numerische Mathematik}, 1(1), 269--271.

	\bibitem{strassen1969}
	Strassen, V.\ (1969).
	Gaussian elimination is not optimal.
	\textit{Numerische Mathematik}, 13(4), 354--356.

	\bibitem{coppersmith1990}
	Coppersmith, D.\ \& Winograd, S.\ (1990).
	Matrix multiplication via arithmetic progressions.
	\textit{Journal of Symbolic Computation}, 9(3), 251--280.

	\bibitem{hopfield1982}
	Hopfield, J.\,J.\ (1982).
	Neural networks and physical systems with emergent collective computational abilities.
	\textit{Proceedings of the National Academy of Sciences}, 79(8), 2554--2558.

	\bibitem{hodgkin1952}
	Hodgkin, A.\,L.\ \& Huxley, A.\,F.\ (1952).
	A quantitative description of membrane current and its application to conduction and excitation in nerve.
	\textit{Journal of Physiology}, 117(4), 500--544.

	\bibitem{galvani1791}
	Galvani, L.\ (1791).
	\textit{De viribus electricitatis in motu musculari commentarius}.
	Typographia Instituti Scientiarium, Bologna.

	\bibitem{hebb1949}
	Hebb, D.\,O.\ (1949).
	\textit{The Organization of Behavior: A Neuropsychological Theory}.
	Wiley, New York.

	\bibitem{bliss1973}
	Bliss, T.\,V.\,P.\ \& Lømo, T.\ (1973).
	Long-lasting potentiation of synaptic transmission in the dentate area of the anaesthetized rabbit following stimulation of the perforant path.
	\textit{Journal of Physiology}, 232(2), 331--356.

	\bibitem{delcastillo1954}
	Del Castillo, J.\ \& Katz, B.\ (1954).
	Quantal components of the end-plate potential.
	\textit{Journal of Physiology}, 124(3), 560--573.

	\bibitem{markram1997}
	Markram, H., Lübke, J., Frotscher, M.\ \& Sakmann, B.\ (1997).
	Regulation of synaptic efficacy by coincidence of postsynaptic APs and EPSPs.
	\textit{Science}, 275(5297), 213--215.

	\bibitem{goldmanrakic1995}
	Goldman-Rakic, P.\,S.\ (1995).
	Cellular basis of working memory.
	\textit{Neuron}, 14(3), 477--485.

	\bibitem{radley2004}
	Radley, J.\,J., Sisti, H.\,M., Hao, J., Rocher, A.\,B., McCall, T., Hof, P.\,R., McEwen, B.\,S.\ \& Morrison, J.\,H.\ (2004).
	Chronic behavioral stress induces apical dendritic reorganization in pyramidal neurons of the medial prefrontal cortex.
	\textit{Cerebral Cortex}, 14(7), 711--717.

	\bibitem{lupien2009}
	Lupien, S.\,J., McEwen, B.\,S., Gunnar, M.\,R.\ \& Heim, C.\ (2009).
	Effects of stress throughout the lifespan on the brain, behaviour and cognition.
	\textit{Nature Reviews Neuroscience}, 10(6), 434--445.

	\bibitem{mcewan2007}
	McEwen, B.\,S.\ (2007).
	Physiology and neurobiology of stress and adaptation: central role of the brain.
	\textit{Physiological Reviews}, 87(3), 873--904.

	\bibitem{pariante2008}
	Pariante, C.\,M.\ \& Lightman, S.\,L.\ (2008).
	The HPA axis in major depression: classical theories and new developments.
	\textit{Trends in Neurosciences}, 31(9), 464--468.

	\bibitem{fox2005}
	Fox, M.\,D., Snyder, A.\,Z., Vincent, J.\,L., Corbetta, M., Van Essen, D.\,C.\ \& Raichle, M.\,E.\ (2005).
	The human brain is intrinsically organized into dynamic, anticorrelated functional networks.
	\textit{Proceedings of the National Academy of Sciences}, 102(27), 9673--9678.

	\bibitem{buckner2008}
	Buckner, R.\,L., Andrews-Hanna, J.\,R.\ \& Schacter, D.\,L.\ (2008).
	The brain's default network: anatomy, function, and relevance to disease.
	\textit{Annals of the New York Academy of Sciences}, 1124, 1--38.

	\bibitem{buzsaki2004}
	Buzsáki, G.\ \& Draguhn, A.\ (2004).
	Neuronal oscillations in cortical networks.
	\textit{Science}, 304(5679), 1926--1929.

	\bibitem{fries2005}
	Fries, P.\ (2005).
	A mechanism for cognitive dynamics: neuronal communication through neuronal coherence.
	\textit{Trends in Cognitive Sciences}, 9(10), 474--480.

	\bibitem{singer1999}
	Singer, W.\ (1999).
	Neuronal synchrony: a versatile code for the definition of relations?
	\textit{Neuron}, 24(1), 49--65.

	\bibitem{varela2001}
	Varela, F., Lachaux, J.\,P., Rodriguez, E.\ \& Martinerie, J.\ (2001).
	The brainweb: phase synchronization and large-scale integration.
	\textit{Nature Reviews Neuroscience}, 2(4), 229--239.

	\bibitem{canolty2006}
	Canolty, R.\,T., Edwards, E., Dalal, S.\,S., Soltani, M., Nagarajan, S.\,S., Kirsch, H.\,E., Berger, M.\,S., Barbaro, N.\,M.\ \& Knight, R.\,T.\ (2006).
	High gamma power is phase-locked to theta oscillations in human neocortex.
	\textit{Science}, 313(5793), 1626--1628.

	\bibitem{mormann2005}
	Mormann, F., Fell, J., Axmacher, N., Weber, B., Lehnertz, K., Elger, C.\,E.\ \& Fernández, G.\ (2005).
	Phase/amplitude reset and theta-gamma interaction in the human medial temporal lobe during a continuous word recognition memory task.
	\textit{Hippocampus}, 15(7), 890--900.

	\bibitem{klimesch1999}
	Klimesch, W.\ (1999).
	EEG alpha and theta oscillations reflect cognitive and memory performance: a review and analysis.
	\textit{Brain Research Reviews}, 29(2--3), 169--195.

	\bibitem{tononi2014}
	Tononi, G.\ \& Cirelli, C.\ (2014).
	Sleep and the price of plasticity: from synaptic and cellular homeostasis to memory consolidation and integration.
	\textit{Neuron}, 81(1), 12--34.

	\bibitem{weaver2004}
	Weaver, I.\,C.\,G., Cervoni, N., Champagne, F.\,A., D'Alessio, A.\,C., Sharma, S., Seckl, J.\,R., Dymov, S., Szyf, M.\ \& Meaney, M.\,J.\ (2004).
	Epigenetic programming by maternal behavior.
	\textit{Nature Neuroscience}, 7(8), 847--854.

	\bibitem{crick1990}
	Crick, F.\ \& Koch, C.\ (1990).
	Towards a neurobiological theory of consciousness.
	\textit{Seminars in the Neurosciences}, 2, 263--275.

	\bibitem{dehaene2011}
	Dehaene, S.\ \& Changeux, J.\,P.\ (2011).
	Experimental and theoretical approaches to conscious processing.
	\textit{Neuron}, 70(2), 200--227.

	\bibitem{friston2010}
	Friston, K.\ (2010).
	The free-energy principle: a unified brain theory?
	\textit{Nature Reviews Neuroscience}, 11(2), 127--138.

	\bibitem{miller1956}
	Miller, G.\,A.\ (1956).
	The magical number seven, plus or minus two: some limits on our capacity for processing information.
	\textit{Psychological Review}, 63(2), 81--97.

	\bibitem{raichle2006}
	Raichle, M.\,E.\ \& Mintun, M.\,A.\ (2006).
	Brain work and brain imaging.
	\textit{Annual Review of Neuroscience}, 29, 449--476.

	\bibitem{fricker2007}
	Fricker, M.\ (2007).
	\textit{Epistemic Injustice: Power and the Ethics of Knowing}.
	Oxford University Press, Oxford.

	\end{thebibliography}

\end{document}
